\documentclass[11pt]{amsart}
\usepackage[T1]{fontenc}
\usepackage{lmodern,amsmath,amssymb,amsthm,mathtools,microtype,booktabs,array,longtable}
\usepackage[margin=1.06in]{geometry}
\usepackage[hidelinks]{hyperref}
\emergencystretch=2em
\numberwithin{equation}{section}
\newtheorem{theorem}{Theorem}[section]
\newtheorem{lemma}[theorem]{Lemma}
\newtheorem{proposition}[theorem]{Proposition}
\newtheorem{corollary}[theorem]{Corollary}
\theoremstyle{definition}\newtheorem{definition}[theorem]{Definition}
\theoremstyle{remark}\newtheorem{remark}[theorem]{Remark}
\newcommand{\Q}{\mathbb Q}\newcommand{\R}{\mathbb R}\newcommand{\C}{\mathbb C}\newcommand{\N}{\mathbb N}\newcommand{\Z}{\mathbb Z}\newcommand{\E}{\mathbb E}
\newcommand{\ip}[2]{\langle #1,#2\rangle}\newcommand{\norm}[1]{\lVert #1\rVert}
\newcommand{\epsc}{\varepsilon_{\mathrm c}}
\DeclareMathOperator{\osc}{osc}\DeclareMathOperator{\spanop}{span}\DeclareMathOperator{\End}{End}\DeclareMathOperator{\diag}{diag}\DeclareMathOperator{\tr}{tr}\DeclareMathOperator{\rad}{rad}\DeclareMathOperator{\TV}{TV}
\title[Stochastic purification and noise thresholds]{Stochastic Purification and Sharp Noise Thresholds\\for Compact Group Experiments}
\author{Qian Qi}
\date{27 September 2026}
\subjclass[2020]{Primary 93B15, 68Q45; Secondary 22C05, 60J10, 11J13}
\keywords{stochastic realization, compact group, categorical output, Chebyshev radius, nonuniform width, Diophantine approximation}
\begin{document}
\begin{abstract}
We prove that every stationary finite stochastic realization of a continuous compact-group interface can be replaced, without increasing its number of labels or its uniform error, by one with permutation transitions. The initialization may remain random. No return word or finiteness assumption on the component group is required. This gives a finite-permutation characterization of the unrestricted stationary minimum and an exact finite-quotient criterion for its existence; a cyclic example separates five stochastic labels from six deterministic quotient labels. For horizon-dependent clocked realizations, the constrained component-image Chebyshev radius is the sharp threshold under executable returns. Every nonempty exact-return language is eventually periodic, and we identify the threshold and narrow-cut occupation obstruction in each length phase. For groups with infinitely many components we prove the strict subcritical obstruction and supercritical realization theorem, characterize stationary boundary attainment, and settle the zero-error clocked boundary by finite observable quotients. An explicit profinite experiment has linear exact width and bounded width at every positive error. These statements apply to entire categorical laws and compact-convex outputs. For a specified Diophantine rotation alphabet we retain the matching polynomial width law throughout the subcritical interval, and extend it to unit signal amplitude at every positive subcritical error. Rational orthogonal, polynomial categorical data admit a terminating stationary minimization procedure. The resource is nonuniform stochastic label width; arbitrary real descriptions and exact sampling are not charged.
\end{abstract}
\maketitle
\section{Introduction}
A finite stochastic register has a continuum of probability distributions. It is therefore not immediate that reversibility of a numerical experiment forces a finite-state implementation to be deterministic. The implementation need not respect the group relations: a physically trivial word may act nontrivially on its hidden register. Approximate agreement makes this distinction especially important. Our first result shows that, for one stationary implementation valid at every word, the transition randomness can nevertheless be removed without a loss in width or accuracy. The proof retains the physical group coordinate throughout the reduction.

Throughout the stationary part of the paper, $H$ is an arbitrary compact Hausdorff group and $\mathcal A$ is a finite nonempty alphabet with letters $u_a\in H$ generating a dense subgroup. We use the execution convention
\[
             u_{a_1\cdots a_n}=u_{a_n}\cdots u_{a_1}.
\]
There are finite nonempty sets of seeds $\mathcal S$ and queries $\mathcal J$, and continuous maps $F_{s,j}:H\to Y_j$. Each $Y_j$ is a nonempty compact convex subset of a finite-dimensional real normed space. A categorical output is the whole vector in a probability simplex, with total variation as its norm. Distinct incompatible measurements are separate queries, not a postulated joint distribution.

\begin{definition}[Stationary realization]\label{def:stationary55}
A realization on $k$ labels consists of initialization rows $\alpha_s\in\Delta_k$, row-stochastic matrices $T_a$ of order $k$, and decoder rows $D_j(i)\in Y_j$. All are independent of the word length and clock. Its error is
\[
 \sup_{s,j,w\in\mathcal A^*}
 \norm{\alpha_sT_{a_1}\cdots T_{a_n}D_j-F_{s,j}(u_w)}_j.
\]
The empty word is included. Write $M_\varepsilon$ for the least such $k$ at error at most $\varepsilon$, or $\infty$ if there is none. A permutation realization has permutation matrices for all commands; its initialization need not be a point mass.
\end{definition}
All advertised labels count, even if some are unused. A decoder vector is not an additional sampled-answer register. When an actual finite answer register is charged, its size is included separately as in Definition~\ref{def:machine54}.

Theorem~\ref{thm:purification55} proves that the minimum $M_\varepsilon$ is unchanged by restricting to permutation realizations. Theorem~\ref{thm:stationaryminimum55} expresses it by finitely many permutation choices at each $k$, followed by a compact feasibility problem. In particular, the transition matrices need not be searched over a continuum. Random initialization remains essential: Proposition~\ref{prop:C655} has unrestricted stationary minimum five and deterministic component-quotient minimum six.

For an open normal subgroup $U$ of $H$ put
\begin{equation}\label{eq:quotientradius55}
 r(U)=\max_{s,j,g\in H}\rad_{Y_j}\bigl(F_{s,j}(gU)\bigr),
 \qquad
 \rad_Y(A)=\min_{y\in Y}\sup_{x\in A}\norm{x-y}.
\end{equation}
These maxima and minima are attained: continuous images of compact cosets are compact, and the radius depends continuously on their translate. The finite-quotient criterion, Theorem~\ref{thm:finitequotient55}, says precisely
\[
          M_\varepsilon<\infty
          \quad\Longleftrightarrow\quad
          r(U)\le\varepsilon\text{ for some open normal }U.
\]
This is an equality-inclusive criterion, valid without an executable return. If $H/H^\circ$ is finite, it reduces to the largest constrained radius of a connected-component image.

A different resource is the least width $W_{N,\varepsilon}$ when all command rows may depend on a free clock and the entire machine may be redesigned at horizon $N$. The stationary purification theorem does not identify this resource with $M_\varepsilon$. The distinction is necessary: a single irrational-rotation letter admits a one-label clocked realization at each horizon. For clocked realizations we prove a converse by localized conditional-moment contraction. With cofinite exact returns and finitely many components, it gives the sharp component-radius theorem, including boundary attainment and an occupation bound for all cuts of any fixed width. Section~\ref{sec:phases55} replaces cofiniteness by the existence of a nonempty exact return, with the necessary refinement by length phase. Theorem~\ref{thm:profinite55} treats arbitrary component groups on both strict sides of the radius, while Theorem~\ref{thm:zero55} identifies the zero-error boundary with finite observable quotients. A concrete $2$-adic interface has $W_{N,0}=\Theta(N)$ and an all-width occupation bound, but bounded width at every positive error (Theorem~\ref{thm:adiclinear55}).

The quantitative theory has separate arithmetic hypotheses. For $s$ planar rotations with a dual badly approximable angle vector, the width is $\Theta(N^{s/(2s+1)})$ at every fixed error below half the signal amplitude when $0<\rho<1$. At $\rho=1$ the same law holds at each positive subcritical error, by Corollary~\ref{cor:unitnoise55}; the exact endpoint is not inferred by a limiting argument. Neither rational matrix entries nor an arbitrary enlarged alphabet implies the hypotheses of the matching upper bound.

The matrix-group structure used in purification is classical; in particular, compact groups of nonnegative matrices and their permutation structure were studied by Brown, Schwarz and Flor \cite{Flor}. Our reduction applies that structure to the joint closure of physical products and stochastic rows and preserves every seed--query error constraint. We give the necessary argument rather than assume that group relations already hold on hidden states. The remaining proofs use Haar projection, representative functions, constrained radius localization, executable word packing, and polygon enclosure. Section~\ref{sec:comparison54} separates these ingredients from the realization statements established here.

\section{Stationary stochastic purification}\label{sec:purification55}
Write $\mathsf{Stoch}(k)$ for the compact semigroup of row-stochastic matrices of order $k$. The identity of a group inside this semigroup can be a singular idempotent; it must not be confused with the identity matrix.

\begin{lemma}[An idempotent over the physical identity]\label{lem:idempotent55}
Let $\mathcal T$ be a compact subsemigroup of $H\times\mathsf{Stoch}(k)$ whose projection onto $H$ is surjective. There is an element $e=(1,E)\in\mathcal T$ with $E^2=E$ and rank equal to the minimum matrix rank in $\mathcal T$. The monoid $e\mathcal T e$ is a compact group with identity $e$ and still projects onto $H$.
\end{lemma}
\begin{proof}
Choose $(h,T)$ of minimum rank $r$. The powers of $T$ are bounded, since they are stochastic. Thus its complex eigenvalues have modulus at most one and all Jordan blocks on the unit circle have size one. The closure of positive powers of an element of a compact group is a group: arbitrarily large positive powers approach its identity, and the preceding powers approach its inverse. Apply this observation in the product of $H$ and the finitely many unit circles corresponding to the peripheral eigenvalues of $T$. There is a net of integers tending to infinity for which $h^n\to1$ and all peripheral eigenvalues raised to the $n$th power tend to one. Along it $T^n$ tends to the spectral projection $E$ onto the peripheral subspace. This is a real stochastic idempotent. Its rank is at most $r$ and at least $r$ by minimality, so it is $r$. Nets avoid any metrizability assumption on $H$.

Every matrix $B$ in $e\mathcal T e$ satisfies $B=EBE$ and has rank $r$. Its action on the row space of $E$ is consequently invertible. Restriction embeds $e\mathcal T e$ as a compact submonoid of $H\times\mathrm{GL}(r,\R)$, with identity $(1,I_r)$; injectivity follows from $B=EBE$. A compact submonoid of a topological group is a subgroup, by the same positive-power return argument. The inverse images of its inverses are still stochastic. Finally, sandwiching an element with physical coordinate $h$ between $e$'s leaves that coordinate unchanged. Surjectivity follows.
\end{proof}

\begin{lemma}[The recurrent simplex]\label{lem:simplex55}
If $E$ is a stochastic idempotent of rank $r$, then
\[
 \Delta_k E=\{p\in\Delta_k:pE=p\}
\]
is a simplex with exactly $r$ vertices. Every group of stochastic matrices with identity $E$ permutes these vertices.
\end{lemma}
\begin{proof}
Regard $E$ as a transition matrix. Its closed communicating classes each have a unique stationary probability row supported on that class. Every stationary row is a convex combination of these rows: mass on a transient class is zero, and stationarity within a closed irreducible class determines a one-dimensional space. For completeness, transient mass vanishes because each transient state has a positive probability of reaching a closed class in a bounded number of steps; a common bound over the finitely many states makes transient mass decrease by a fixed factor unless it is zero. The closed-class rows have disjoint supports, hence are linearly independent.

Each row of $E$ is stationary because $E^2=E$. Conversely, every stationary row equals its product with $E$. The span of the closed-class stationary rows is therefore precisely the row space of $E$, so there are $r$ classes and $r$ simplex vertices. If $B$ has a stochastic inverse relative to $E$, multiplication by $B$ and by this inverse are inverse affine maps of this simplex. Affine bijections permute its vertices.
\end{proof}

\begin{theorem}[Purification without a width or error loss]\label{thm:purification55}
For every compact-group interface in Definition~\ref{def:stationary55}, a stationary $k$-label realization of error at most $\varepsilon$ admits a permutation realization of error at most $\varepsilon$ on at most $k$ labels. There is no return-word hypothesis and no restriction on the number of connected components of $H$.
\end{theorem}
\begin{proof}
For the given machine form
\[
 \mathcal T=\overline{\{(u_w^{-1},T_w):w\in\mathcal A^*\}}
       \subset H\times\mathsf{Stoch}(k).
\]
The inverse on the physical coordinate is important: concatenation satisfies $u_{vw}^{-1}=u_v^{-1}u_w^{-1}$ and $T_{vw}=T_vT_w$. Thus $\mathcal T$ is a compact semigroup. Positive-word density makes its projection onto $H$ surjective. By continuity, its every element $(h,B)$ satisfies
\begin{equation}\label{eq:closederror55}
             \norm{\alpha_sBD_j-F_{s,j}(h^{-1})}_j\le\varepsilon
             \quad(s\in\mathcal S,\ j\in\mathcal J).
\end{equation}

Apply Lemma~\ref{lem:idempotent55} and denote the vertices of $\Delta_k E$ by $\nu_1,\ldots,\nu_r$. For each letter the element
\[
        e(u_a^{-1},T_a)e=(u_a^{-1},ET_aE)
\]
belongs to the group $e\mathcal T e$. By Lemma~\ref{lem:simplex55}, $B_a=ET_aE$ permutes the $\nu_i$. Use that permutation as the new command. Write uniquely
\[
          \alpha_sE=\sum_{i=1}^r\beta_s(i)\nu_i,
          \qquad \widetilde D_j(i)=\nu_iD_j\in Y_j.
\]
The last inclusion uses convexity of the legal output set. For a nonempty word the output of the resulting $r$-label machine is $\alpha_s B_{a_1}\cdots B_{a_n}D_j$. For the empty word it is $\alpha_sED_j$. In both cases the corresponding joint element lies in $\mathcal T$ and has physical coordinate $u_w^{-1}$. Equation~\eqref{eq:closederror55} proves the error bound. Since $r\le k$, the width has not increased.
\end{proof}
The new machine need not reproduce the old machine's erroneous output word by word; it remains in the same closed error balls about the target. Nor do we assert $T_{w}=I$ on a physical return, invertibility of the original $T_a$, or deterministic initialization. These three incorrect shortcuts would invalidate the reduction.

\subsection{Finite quotients and the unrestricted stationary minimum}
For a tuple of permutation matrices $\sigma=(\Pi_a)_{a\in\mathcal A}$ of order $k$, set
\[
 \mathcal L_\sigma=
 \overline{\{(u_w^{-1},\Pi_w):w\in\mathcal A^*\}}
       \subset H\times\mathfrak S_k.
\]
This is a compact subgroup. Its definition includes the possible discrepancy between physical relations and hidden-state relations.

\begin{theorem}[An exact stationary minimum]\label{thm:stationaryminimum55}
The value $M_\varepsilon$ is the least positive integer $k$ for which some tuple $\sigma$ of permutations of order $k$, some $\beta_s\in\Delta_k$, and some $d_j(i)\in Y_j$ satisfy
\begin{equation}\label{eq:permutationfeasibility55}
 \norm{\beta_s\Pi d_j-F_{s,j}(h^{-1})}_j\le\varepsilon
 \quad\text{for all }(h,\Pi)\in\mathcal L_\sigma\text{ and all }s,j.
\end{equation}
For fixed $k$ there are at most $(k!)^{|\mathcal A|}$ discrete choices, and the remaining feasible set is compact. The optimal stationary error at any fixed width is attained.
\end{theorem}
\begin{proof}
Necessity follows from Theorem~\ref{thm:purification55}, padding by unused labels if its output has fewer than $k$ states, and then taking the closure of word products. Sufficiency is the restriction of \eqref{eq:permutationfeasibility55} to words. The probabilities and legal decoders lie in a finite product of compact sets. Each displayed inequality defines a closed subset of that product. The worst error is continuous in these parameters, uniformly in the compact group variable, so its minimum exists. There are finitely many permutation tuples.
\end{proof}

\begin{theorem}[Finite observable quotient criterion]\label{thm:finitequotient55}
For arbitrary compact $H$ and the stated continuous interface,
\[
 M_\varepsilon<\infty
 \quad\Longleftrightarrow\quad
 r(U)\le\varepsilon\text{ for some open normal subgroup }U\le H.
\]
If such $U$ exists, $|\mathcal S|[H:U]$ deterministic-initialized permutation labels suffice. Conversely, a stationary $k$-label realization supplies such a $U$ with $[H:U]\le k!$.
\end{theorem}
\begin{proof}
Purify a stationary realization and use its joint group $\mathcal L_\sigma$. Put
\[
                  U=\{h\in H:(h,I)\in\mathcal L_\sigma\}.
\]
This is a closed normal subgroup of $H$: conjugate $(h,I)$ by an element above any $g\in H$, using surjectivity of the first projection. For a fixed permutation $\Pi$, a nonempty fiber in $\mathcal L_\sigma$ is a coset $hU$. These at most $k!$ fibers cover $H$, so $U$ has finite index. A closed subgroup of finite index in a compact group is open, since its complement is a finite union of closed cosets.

For each $s,j$ and each fiber, the single legal vector $\beta_s\Pi d_j$ is within $\varepsilon$ of every $F_{s,j}(h^{-1})$ in that fiber. Normality identifies the inverse fiber with a coset of $U$. Hence $r(U)\le\varepsilon$. Conversely, store $(s,gU)$, update the coset by left multiplication by $u_a$, and decode a minimizing center of $F_{s,j}(gU)$. Compactness supplies the centers. This gives the asserted finite realization.
\end{proof}
For finitely many connected components, every open subgroup contains $K=H^\circ$, and $K$ itself is open normal. Therefore $M_\varepsilon<\infty$ holds exactly when
\[
 \varepsilon\ge\max_{s,j,c\in H/K}\rad_{Y_j}(F_{s,j}(c)),
\]
\emph{without} a return assumption. Executable returns enter the stronger clocked converse, not this stationary statement. The index bound $k!$ is an existence bound on a deterministic quotient, not a claim that the stationary minimum equals its index.

\begin{proposition}[Random initialization strictly saves states]\label{prop:C655}
Let $H=\Z/6\Z$, with an identity letter and a generator, one seed, and binary success probability
\[
 f(n)=\tfrac12+\tfrac14(-1)^n+\tfrac14\cos(2\pi n/3).
\]
At zero error the least deterministic component-quotient width is six, whereas the unrestricted stationary stochastic minimum is five.
\end{proposition}
\begin{proof}
The values at $n=0,\ldots,5$ are
\[
                  (1,\tfrac18,\tfrac58,\tfrac12,\tfrac58,\tfrac18).
\]
This sequence has least period six. A deterministic quotient of the cyclic component action preserving these outputs must therefore have six states. For a five-state realization take the disjoint union of a two-cycle and a three-cycle, initialize with mass $1/2$ at the zeroth vertex of each, and use decoder rows $(1,0)$ on the two-cycle and $(1,1/4,1/4)$ on the three-cycle. The identity command fixes all five labels. Direct averaging gives exactly $f(n)$, with legal binary probabilities.

If a stationary realization on at most four labels existed, purification would give one whose generator is a permutation of at most four letters. Its order is one of $1,2,3,4$, as follows by listing the possible cycle lengths. Every resulting numerical sequence has a period dividing that order, contradicting the least period six. Thus five is the unrestricted minimum.
\end{proof}

\section{The clocked component-radius threshold}\label{sec:intro54}
A finite stochastic register can carry infinitely many probability distributions. Thus an infinite numerical orbit is not, by itself, an obstruction to finite positive realization. A further difficulty arises when the realization may be redesigned at every horizon and its transitions may depend on a free external clock. The obstruction established here survives both freedoms. It is measured by the radius of an observable image inside its allowed output set.

For the clocked theorem in this section, let $H$ be a compact Hausdorff group with finitely many connected components. Put $K=H^\circ$ and $C=H/K$. Then $K$ is an open normal subgroup and $C$ is a finite group. Fix a finite alphabet $\mathcal A$ with elements $u_a\in H$ whose generated group is dense in $H$. Our convention is
\[
 u_{a_1\cdots a_n}=u_{a_n}\cdots u_{a_1}.
\]
The identity $e$ need not be an alphabet element. Instead assume the following executable return condition:
\begin{equation}\label{eq:returns54}
 \text{there is an integer }g\ge0\text{ such that, for every }n\ge g,
 \text{ some word of length }n\text{ has product }e.
\end{equation}
An identity letter gives $g=0$. The condition concerns physical products; the stochastic rows on a return word may perform arbitrary processing.

For nonempty finite seed and query sets $\mathcal S,\mathcal J$, let
\[
 F_{s,j}:H\longrightarrow Y_j
\]
be continuous, where $Y_j$ is a nonempty compact convex subset of a finite-dimensional real normed space $(V_j,\norm{\cdot}_j)$. A query is selected only after the command word. A prescribed joint law of finitely many outputs is treated as one query, not as a collection of marginal constraints. In particular, a categorical interface has
\[
 Y_j=\Delta_{M_j}=\{p\in\R^{M_j}:p_i\ge0,\ \textstyle\sum_i p_i=1\},
 \qquad \norm{x}_j=\tfrac12\sum_i|x_i|.
\]
Then $\norm{p-q}_j$ is total variation. Binary means correspond to $Y=[-1,1]$ with $\norm{x}=|x|/2$.

\begin{definition}[Clocked realization]\label{def:machine54}
At horizon $N$ choose finite nonempty registers $S_t$, $0\le t\le N$, initialization rows $E_s$, row-stochastic command matrices $T_{t,a}:S_{t-1}\to S_t$, and terminal decoder maps $D_j:S_N\to Y_j$. The realized output vector on $w=a_1\cdots a_N$ is
\[
 E_sT_{1,a_1}\cdots T_{N,a_N}D_j.
\]
Its error is the maximum, over all seeds, length-$N$ words, and queries, of its distance from $F_{s,j}(u_w)$. The command width is $\max_t|S_t|$, counting all declared labels, including unused padding labels. Write $W_{N,\varepsilon}$ for its least possible value at error at most $\varepsilon$.
\end{definition}
All retained private information, including randomness kept for future use, is charged to the current label. The rows may depend on $t$ and $N$, but not on an unrecorded input history or on the future query. The external clock and horizon, arbitrary real table descriptions, their construction and lookup, real arithmetic, and exact sampling are uncharged. This is a numerical realization invariant, not uniform workspace or program size. For a finite categorical interface, charging the terminal answer register changes the peak to $\max(W_{N,\varepsilon},\max_j M_j)$. Thus all boundedness and growth results are unchanged; the convention that also charges a sampled binary-answer register has a minimum of two labels. A decoder vector by itself is not an additional sampled register.

For $A\subset Y_j$ define its constrained radius by
\[
 \rad_{Y_j}(A)=\min_{y\in Y_j}\sup_{x\in A}\norm{x-y}_j.
\]
The minimum is attained when $A$ is compact. The critical error is
\begin{equation}\label{eq:radius54}
 \epsc=\max_{s,j,c\in C}\rad_{Y_j}\bigl(F_{s,j}(c)\bigr).
\end{equation}
Here a component is identified with its coset in $H$; its image is compact by continuity. The requirement that the center belong to $Y_j$ is essential for categorical outputs.

\begin{theorem}[Intrinsic radius classification]\label{thm:radius54}
Under \eqref{eq:returns54}, for every $\varepsilon\ge0$ the following are equivalent:
\begin{enumerate}
\item $\sup_N W_{N,\varepsilon}<\infty$;
\item one stationary finite-state realization with permutation command updates has error at most $\varepsilon$ at every horizon;
\item $\varepsilon\ge\epsc$.
\end{enumerate}
At the boundary, at most $|\mathcal S||C|$ command labels suffice. If $\varepsilon<\epsc$, then for every integer $k\ge1$ there is an integer $L_k<\infty$ such that every error-$\varepsilon$ realization satisfies
\begin{equation}\label{eq:occupation54}
 \#\{t\in\{1,\ldots,N\}:|S_t|\le k\}\le L_k.
\end{equation}
All intervening widths are unrestricted. In particular, $W_{N,\varepsilon}\to\infty$.
\end{theorem}
Neither a positive lower bound on state masses nor a reachable-rank, observability, or conditioning assumption occurs in the theorem. The upper realization is deterministic until its terminal decoder. It need not have the minimum width among all stochastic realizations. Theorem~\ref{thm:stationaryminimum55}, rather than a deterministic component partition, characterizes that stationary minimum.

For binary means, \eqref{eq:radius54} reduces to one quarter of the largest componentwise oscillation. For a genuine categorical law it is generally not half its diameter; an analytic three-outcome example below has critical total-variation error $2/3$, whereas every individual event has binary threshold $1/2$. The proof uses finitely many localized vector tests rather than a separate midpoint for each coordinate.

The finite-component assumption does not assert that $H$ is a Lie group. For example, take $H=\prod_{n\ge1}\R/\Z$, choose an element whose every finite coordinate tuple together with $1$ is rationally independent, and include this element and $e$ as commands. The positive orbit is dense and $H$ is connected but has no faithful finite-dimensional representation. The theorem applies to every finite continuous interface on this group. With infinitely many components the component construction need not be finite; Section~\ref{sec:profinite55} gives the finite-observable-quotient criterion and the strict-side clocked theorem in that setting.

There is also a quantitative counterpart. For planar rotations with a dual badly approximable $s$-tuple of angles and signal amplitude $0<\rho<1$, Theorem~\ref{thm:matched54} proves
\[
 W_{N,\varepsilon}=\Theta\bigl(N^{s/(2s+1)}\bigr)
 \quad\text{for every fixed }0\le\varepsilon<\rho/2.
\]
The upper realization is exact even when $\varepsilon>0$. The new lower estimate reaches the entire subcritical interval, by combining harmonic localization with the arithmetic orbit-packing mechanism. This quantitative theorem and the radius classification have distinct hypotheses: algebraic matrix entries alone do not imply the badly approximable angle condition.

\subsection{Executable tests}
Let $\mathcal P=\{u_w:w\in\mathcal A^*\}$. We use standard Haar averaging and the density of matrix coefficients of finite-dimensional unitary representations, in realified form \cite{Folland}. No quantitative mixing assumption for a preassigned random walk is imposed.

\begin{lemma}[Positive words and common lengths]\label{lem:positive54}
The semigroup $\mathcal P$ is dense in $H$, and $\mathcal P\cap K$ is dense in $K$. Every finite subset of $\mathcal P\cap K$ has executable representatives of a common positive length.
\end{lemma}
\begin{proof}
For an element $u$ of a compact group, positive powers return arbitrarily close to $e$ with exponent at least two. Indeed cover the compact closure of its powers by finitely many translates of a sufficiently small neighborhood. Among sufficiently many powers, three lie in one member of the cover; the first and last exponents differ by at least two, and their quotient is arbitrarily close to $e$. This argument also covers finite-order elements by taking a sufficiently large multiple of the order. If $u^r$ approaches $e$ with $r\ge2$, then $u^{r-1}=u^ru^{-1}$ approaches $u^{-1}$ and is still a positive power. Thus $u^{-1}$ is in the closure of positive powers. The closed semigroup $\overline{\mathcal P}$ therefore contains the group generated by the letters, and hence equals $H$. Since $H$ has finitely many components, $K$ is open, and a relatively open subset of $K$ is open in $H$; density gives the second assertion. For finitely many representative lengths $\ell_i$, choose $B\ge\max_i\ell_i+g$, $B\ge1$, and append return words of lengths $B-\ell_i$. The products are unchanged.
\end{proof}

\begin{lemma}[Finite executable gap]\label{lem:gap54}
Let $R:K\to O(E)$ be a continuous representation on a nonzero finite-dimensional Euclidean space, with no nonzero fixed vector. For every $k\ge1$ there are $B\ge1$, a probability law $\nu$ on finitely many executable length-$B$ words with products in $K$, and $d\in(0,1)$ such that
\begin{equation}\label{eq:gap54}
 \E_{w\sim\nu}\max_{1\le i\le k}\ip{v_i}{R(u_w)u}\le1-d
\end{equation}
for all unit vectors $u,v_1,\ldots,v_k\in E$.
\end{lemma}
\begin{proof}
On the compact product of unit spheres consider the minimum
\[
 \gamma=\min_{u,v_1,\ldots,v_k}\int_K
 \left(1-\max_i\ip{v_i}{R(h)u}\right)\,d\mu(h).
\]
Repeated centers are permitted. The integrand is nonnegative. If a minimizing tuple had integral zero, full support of Haar measure would give $R(h)u\in\{v_1,\ldots,v_k\}$ for every $h$. The orbit of $u$ is connected, hence a singleton, contradicting the fixed-vector hypothesis. Thus $\gamma>0$.

Partition the compact image $R(K)$ into finitely many measurable sets of small diameter in the Euclidean operator norm, and replace their Haar masses by atoms at nearby executable products in $K$. Cells of zero Haar mass can be discarded. The function $A\mapsto\max_i\ip{v_i}{Au}$ is uniformly $1$-Lipschitz in that norm. An approximation radius less than $\gamma/2$ therefore preserves a positive defect uniformly in the unit vectors. Lemma~\ref{lem:positive54} gives a common length. Decreasing the defect to a number in $(0,1)$ proves the assertion.
\end{proof}

\begin{lemma}[Conditional-centroid contraction]\label{lem:centroid54}
Let $Z\in E$ be a bounded random physical vector and $S$ the current register. Apply an independent word $w\sim\nu$ from Lemma~\ref{lem:gap54}, and let the actual block kernel generate a terminal state $S'$ having at most $k$ labels. Then
\begin{equation}\label{eq:centroid54}
 \E\norm{\E[R(u_w)Z\mid S']}\le(1-d)\E\norm{\E[Z\mid S]}.
\end{equation}
A fixed intervening return word cannot increase the same conditional norm, regardless of its intermediate or terminal widths.
\end{lemma}
\begin{proof}
Put $z_s=\E[Z\mid S=s]$ and write $T_w(s,i)$ for the actual composite row of the block. Since this row depends on the past only through $S$,
\[
 \Pr(S'=i)\E[R(u_w)Z\mid S'=i]
 =\sum_s\Pr(S=s)\E_w T_w(s,i)R(u_w)z_s.
\]
Choose $v_i$ in the direction of each nonzero terminal centroid, arbitrarily otherwise. Pair, sum in $i$, and use stochasticity to bound the sum by
\[
 \sum_s\Pr(S=s)\E_w\max_i\ip{v_i}{R(u_w)z_s}
 \le(1-d)\sum_s\Pr(S=s)\norm{z_s}.
\]
Zero-probability states contribute zero. No restriction was placed on registers inside the block. On a fixed return word the physical vector is unchanged at its endpoints; the new state is a stochastic image of the old one, so conditional Jensen gives nonincrease. Independence is used only at complete block boundaries.
\end{proof}

\section{Representative functions and localized radii}\label{sec:local54}
A real representative function on $K$ is a matrix coefficient of a finite-dimensional continuous orthogonal representation, or a finite sum of such coefficients. Their algebra contains the constants, is closed under multiplication by tensor products, and is uniformly dense in $C(K,\R)$ by the Peter--Weyl approximation theorem \cite{Folland}. A finite collection of representative functions can be written on one space
\[
 \mathcal V=\bigoplus_{i=1}^a\End(E_i),\qquad
 \Phi(h)=(R_i(h))_{i=1}^a,
\]
where the scalar product is the sum of Frobenius products. Left multiplication defines an orthogonal representation $R$ on $\mathcal V$ and satisfies $\Phi(gh)=R(g)\Phi(h)$. Adjoin a trivial summand when necessary. Let
\[
 P=\int_KR(h)\,d\mu(h),\quad E=(I-P)\mathcal V,\quad
 Z(h)=(I-P)\Phi(h),\quad Q_0=\norm{Z(e)}.
\]
Haar averaging is the orthogonal projection onto fixed vectors. Consequently $Z(gh)=R(g)Z(h)$, the representation on $E$ has no fixed vectors, and $P\Phi(h)=P\Phi(e)$ for $h\in K$.

Choose and fix a coefficient vector $a_p$ for each scalar representative function $p$, and a linear coefficient map $A_G$ for each vector-valued representative function $G$. Their norms need not be intrinsic; fixing any such choices suffices. For an arbitrary joint distribution of $h\in K$ and a finite state $S$, put $Q=\E\norm{\E[Z(h)\mid S]}$ and $\bar p=\int p\,d\mu$. The fixed/moving decomposition gives
\begin{align}
 \E[p(h)\mid S]&=\bar p+\ip{a_p}{\E[Z(h)\mid S]},\label{eq:calibration54}\\
 |\E p-\bar p|&\le\norm{a_p}Q,\qquad
 \norm{\E G-\bar G}\le\norm{A_G}Q.\label{eq:vectorcal54}
\end{align}
Here the norm of $A_G$ is the operator norm from the Euclidean feature space to the indicated output space. If $D:S\to Y$ and $M_Y=\max_{y\in Y}\norm{y}$, then
\begin{equation}\label{eq:decodercal54}
 \norm{\E[p(h)D(S)]-\bar p\E D(S)}\le M_Y\norm{a_p}Q.
\end{equation}
Indeed condition on $S$ in \eqref{eq:calibration54}, multiply by $D(S)$, and apply the triangle inequality. These are moment estimates, not total-variation convergence of a hidden physical distribution.

\begin{lemma}[Finite localization of a radius]\label{lem:radiuslocal54}
Let $f:K\to Y$ be continuous, with $Y$ compact and convex in a finite-dimensional normed space. Suppose $\rad_Y(f(K))>\varepsilon$. There are a vector-valued representative function $G$, a number $\eta>0$, and nonnegative representative functions $p_1,\ldots,p_a$ of Haar mean one such that
\[
 \norm{G-f}_\infty<\eta,\qquad
 r_*:=\min_{y\in Y}\max_i\norm{\bar{p_iG}-y}>\delta:=\varepsilon+\eta.
\]
\end{lemma}
\begin{proof}
Choose $\varepsilon<t<\rad_Y(f(K))$. For each $y\in Y$ some $h\in K$ has $\norm{f(h)-y}>t$. The corresponding open subsets of $Y$ cover $Y$, so a finite subcover yields $h_1,\ldots,h_a$ with
\[
 \min_{y\in Y}\max_i\norm{f(h_i)-y}>t.
\]
The strict inequality follows also at a minimizing $y$, since the maximum is everywhere greater than $t$ and $Y$ is compact.

Fix a sufficiently small positive $\eta$. Approximate the coordinates of $f$ uniformly by representative functions to obtain $\norm{G-f}_\infty<\eta$. For each $i$, choose a nonzero nonnegative continuous bump supported in a neighborhood where $f$ differs from $f(h_i)$ by less than $\eta$. Such a bump exists by normality of the compact Hausdorff space; it has positive Haar integral by full support. Approximate its square root uniformly by a representative function, square the approximation, and normalize its Haar integral. The resulting nonnegative representative function $p_i$ has mean one, and its weighted $f$-mean can be made arbitrarily close to that of the bump. Thus we may ensure
\[
 \norm{\bar{p_iG}-f(h_i)}<3\eta.
\]
The constrained radius changes by at most the maximum perturbation of the finitely many points. Taking $4\eta<t-\varepsilon$ gives $r_*>\varepsilon+\eta$.
\end{proof}

\begin{proposition}[Residual moment from a vector interface]\label{prop:residual54}
Fix the data of Lemma~\ref{lem:radiuslocal54}. Choose one feature space representing all $p_i$ and $p_iG$, and fix their coefficient choices before defining
\[
 L=\max_i\bigl(\norm{A_{p_iG}}+(M_Y+\delta)\norm{a_{p_i}}\bigr),
 \qquad \Delta=r_*-\delta>0.
\]
Suppose a terminal decoder $D(S)\in Y$ obeys
\[
 \norm{\E[D(S)\mid w]-G(h(w))}\le\delta
\]
for every complete deterministic word in a test law. Then $L>0$, $E\ne\{0\}$, $Q_0>0$, and
\begin{equation}\label{eq:residual54}
                         Q\ge\Delta/L.
\end{equation}
\end{proposition}
\begin{proof}
For $p=p_i\ge0$, wordwise correctness and the triangle inequality imply
\[
 \norm{\E[p(h)D(S)]-\E[p(h)G(h)]}\le\delta\E p(h).
\]
This is a conditional mean error, not a pointwise bound on the random decoder. Since $\bar p_i=1$, calibration gives explicitly
\[
                  \E p_i\le1+\norm{a_{p_i}}Q.
\]
Write $y=\E D(S)$, which belongs to $Y$ by convexity. Equations~\eqref{eq:vectorcal54}--\eqref{eq:decodercal54}, together with $\bar p=1$, give
\[
 \norm{y-\bar{p_iG}}\le\delta+
 \bigl(\norm{A_{p_iG}}+(M_Y+\delta)\norm{a_{p_i}}\bigr)Q.
\]
Maximizing in $i$ yields $r_*\le\delta+LQ$. Hence $LQ\ge\Delta>0$. In particular the moving feature and its initial norm cannot vanish: by equivariance, $Q_0=0$ would imply $Z(h)=0$ for all $h$ and therefore $Q=0$.
\end{proof}

\subsection{Proof of Theorem~\ref{thm:radius54}}
Choose $(s,j,c)$ with component radius greater than $\varepsilon$. By density there is an executable prefix $v$ with $u_v\in c$. Put $b=|v|$ and $f(h)=F_{s,j}(hu_v)$ on $K$. Since $Ku_v=c$, this map has the same constrained radius. Apply Lemma~\ref{lem:radiuslocal54} and Proposition~\ref{prop:residual54}.

The moving feature is nonzero independently of any particular test law. Otherwise all represented functions would be constant on $K$; since each $p_i$ has mean one, this would force $p_i=1$ and $G$ constant. Then all $\bar{p_iG}$ would equal a point at distance less than $\eta$ from $Y$, contradicting $r_*>\delta\ge\eta$.

For a fixed width $k$, apply Lemma~\ref{lem:gap54} to the resulting moving feature. Let $B,d$ be its block parameters and set
\[
 \chi=-\log(1-d)>0,\qquad A=\max(1,LQ_0/\Delta).
\]
Start a test word with $v$; between independent length-$B$ test blocks use fixed return words. The accumulated physical product is $hu_v$. Initially $h=e$, so the conditional feature norm after the prefix equals $Q_0$, regardless of the hidden-state distribution. If $J$ block endpoints have width at most $k$, Lemma~\ref{lem:centroid54} gives
\[
 Q_N\le Q_0(1-d)^J.
\]
The machine is correct on every deterministic word of length $N$. Thus any external law supported on such words inherits the conditional mean bound, with error $\delta=\varepsilon+\eta$ for $G$. Proposition~\ref{prop:residual54} implies
\begin{equation}\label{eq:J54}
                         J\chi\le\log A.
\end{equation}
The law depends on the advertised width profile, not on the machine's private state. When $A=1$, \eqref{eq:J54} excludes even one selected narrow block; no logarithmic loss has occurred.

For clarity we give the spacing count with the full return hypothesis. Discard narrow cuts before $b+B+g$ and after $N-g$. From the remaining cuts greedily choose the first and then the next at distance at least $B+g$. Put an independent block immediately before each selected cut. The initial gap, all intervening gaps, and the final suffix have lengths at least $g$, so \eqref{eq:returns54} fills them with executable returns. At most $b+B+2g-1$ cuts were discarded. Each chosen endpoint accounts for at most $B+g$ eligible cuts. Therefore one may take
\begin{equation}\label{eq:Lk54}
 L_k=\left\lceil b+B+2g-1+(B+g)\frac{\log A}{\chi}\right\rceil.
\end{equation}
The constants $B,d$ depend on the fixed moving representation and $k$; $\chi=-\log(1-d)$ depends on that defect. The localization, prefix length $b$, and $A$ depend only on the alphabet, interface, chosen component, and error. Consequently $L_k$ depends on these fixed data, the return certificate $g$, and $k$, but not on $N$ or the machine. When the eligible interval is empty the same estimate follows directly by counting. This proves occupation without restricting any register inside a block or gap. If the entire peak is at most $k$, then $N\le L_k$, which proves divergence directly. Monotonicity in the horizon is not needed when an identity letter is absent.

For the upper bound, use labels $(s,c)\in\mathcal S\times C$, including any labels not reached at a given horizon. Initialize at $(s,K)$ and update by left multiplication
\[
 (s,c)\longmapsto(s,[u_a]c).
\]
Choose a minimizing center of $F_{s,j}(c)$ in $Y_j$ as decoder at $(s,c)$ for query $j$. These are legal decoder values, and the error is at most \eqref{eq:radius54}. All command updates are permutations of one finite state set, independent of clock and horizon. This proves boundary attainment and the equivalence.\qed

\begin{proposition}[A finite return certificate]\label{prop:return54}
Condition~\eqref{eq:returns54} follows if there are finitely many executable identity-product words whose positive lengths have greatest common divisor one.
\end{proposition}
\begin{proof}
Let the lengths be $\ell_1,\ldots,\ell_a$. Their nonnegative sums modulo $\ell_1$ form the subgroup generated by their residues: a subsemigroup of a finite group is a subgroup. The gcd hypothesis makes this the whole residue group. Choose one nonnegative sum $n_r$ for each residue $r$ modulo $\ell_1$, and put $g=\max_r n_r$. Every $n\ge g$ equals its representative $n_r$ plus a nonnegative multiple of $\ell_1$. Concatenating the corresponding return words proves the claim.
\end{proof}
For example, the alphabet $\{R_\theta,R_{-\theta},R_{2\pi/3}\}$, with $\theta/2\pi$ irrational, contains no identity letter but has return lengths two and three. It satisfies \eqref{eq:returns54} with $g=2$. A single irrational rotation letter does not: there is only one word at each horizon, and a horizon-specific constant decoder can store its numerical answer with one command label. Thus synchronization cannot be silently omitted. Nor is a bounded irreversible semigroup enough: for the identity command $I=1$ and contraction command $a=\lambda\in(0,1)$, the mean $\rho\lambda^{\#a(w)}$ has an exact alive/absorbed two-state realization. The group generated with inverses is not bounded. The present theorem is not a compact-semigroup classification.

\section{Exact-return languages and length phases}\label{sec:phases55}
The existence of cofinite returns is not the only useful synchronization regime. In fact an exact-return language, when nonempty, has a rigid eventual form. Here and in this section a return has positive length.

\begin{lemma}[The return period]\label{lem:returnperiod55}
Suppose
\[
 \mathcal R=\{|w|>0:u_w=1\}\ne\varnothing,
 \qquad d=\gcd\mathcal R.
\]
There is an integer $g_d$ such that every multiple $dn$ with $n\ge g_d$ belongs to $\mathcal R$, and no length outside $d\N$ belongs to $\mathcal R$. In particular, every nonempty exact-return language is eventually one arithmetic progression.
\end{lemma}
\begin{proof}
Concatenation makes $\mathcal R$ additive. A finite subset of its elements already has gcd $d$: start with one length and successively decrease the gcd when necessary; a strictly decreasing chain of its positive divisors is finite. Divide these finitely many lengths by $d$. Their gcd is one, so the residue argument of Proposition~\ref{prop:return54} shows that their nonnegative sums contain every sufficiently large integer. Multiplying by $d$ proves the assertion. The exclusion of other lengths is the definition of the gcd.
\end{proof}

Fix a letter with product $a\in H$, and let $J$ be the closure of the products of words whose lengths are multiples of $d$, including the empty word.
\begin{lemma}[The phase subgroup]\label{lem:phasegroup55}
The group $J$ is an open normal subgroup of $H$. The quotient $H/J$ is cyclic of order dividing $d$, generated by $aJ$, and
\begin{equation}\label{eq:phasereachable55}
 \overline{\{u_w:|w|\equiv r\pmod d\}}=a^rJ
 \qquad(0\le r<d).
\end{equation}
If $H/H^\circ$ is finite, then $J^\circ=H^\circ$.
\end{lemma}
\begin{proof}
The compact closure defining $J$ is a group, by the positive-power argument, and is generated densely by the block alphabet $\mathcal A^d$. Also $a^d\in J$. If $x$ is a product of length $d$, then $axa^{d-1}\in J$, whence $axa^{-1}\in J$. It follows that $aJa^{-1}\subset J$. Applying this inclusion $d$ times and using $a^d\in J$ shows equality. Every letter $b$ has $ba^{d-1}\in J$, so $ba^{-1}\in J$. Thus all letters lie in $Ja$, normalize $J$, and their products of length $n$ lie in $Ja^n$. The finite union $\bigcup_{r=0}^{d-1}Ja^r$ is a closed group containing the letters, hence is $H$. This proves normality, finite index, and openness. Conversely, words consisting of $r$ copies of $a$ followed by arbitrary $d$-blocks have products dense in $Ja^r$, proving \eqref{eq:phasereachable55}. An open subgroup contains the identity component, and its own identity component is therefore the same one.
\end{proof}
The quotient order may be a proper divisor of $d$. The phase index records executable length, not merely the physical component quotient.

Assume now that $H$ has finitely many components and put $K=H^\circ$. For $r\in\{0,\ldots,d-1\}$ define
\begin{equation}\label{eq:phaseradius55}
 \varepsilon_r=\max_{s,j,\ c\in H/K:\ c\subset a^rJ}
                    \rad_{Y_j}(F_{s,j}(c)).
\end{equation}
\begin{theorem}[The phasewise clocked threshold]\label{thm:phases55}
Under the sole synchronization assumption $\mathcal R\ne\varnothing$, for every $r$ and $\varepsilon\ge0$,
\[
 \sup_{N\equiv r\ (\mathrm{mod}\ d)}W_{N,\varepsilon}<\infty
             \quad\Longleftrightarrow\quad
             \varepsilon\ge\varepsilon_r.
\]
Equality is attained by a finite component-permutation realization on the indicated horizons. If $\varepsilon<\varepsilon_r$, then for each $k$ there is a constant $L_{k,r}$, independent of $N$, such that every error-$\varepsilon$ realization at $N\equiv r\pmod d$ satisfies
\[
                 \#\{1\le t\le N:|S_t|\le k\}\le L_{k,r}.
\]
The intervening widths are unrestricted. Consequently the full family is bounded exactly at and above $\max_r\varepsilon_r=\varepsilon_{\mathrm c}$; below it, divergence is asserted phase by phase, not necessarily along every horizon.
\end{theorem}
\begin{proof}
Every product of length $N\equiv r$ lies in $a^rJ$. Store its component and the seed and decode component centers as before. This proves the upper bound, including equality, with at most $|\mathcal S||H/K|$ labels. Centers on other components may be chosen arbitrarily when only this phase is under consideration.

For the converse fix a component $c\subset a^rJ$ and $s,j$ with radius greater than $\varepsilon$. We must count narrow cuts in every residue class, not only at block boundaries of one phase. For each $\ell\in\{0,\ldots,d-1\}$ choose $b_\ell\in\{0,\ldots,d-1\}$ congruent to $r-\ell$ modulo $d$. By \eqref{eq:phasereachable55} and the openness of components there is a fixed word $v_\ell$, of length $p_\ell\equiv\ell\pmod d$, with product in $a^{-b_\ell}c$. Indeed this component is contained in $a^\ell J$. Consider the block experiment on $J$ with alphabet $\mathcal A^d$ and target
\[
        G_\ell(h)=F_{s,j}(a^{b_\ell}h u_{v_\ell}),\qquad h\in J.
\]
It is a legal compact-convex interface. Its image on $K=J^\circ$ is $F_{s,j}(c)$, since $K$ is normal. Lemma~\ref{lem:returnperiod55} gives cofinite exact returns in block length for this alphabet.

Restrict the original machine to words that begin with $v_\ell$, end with $b_\ell$ copies of the chosen letter, and have $M=(N-p_\ell-b_\ell)/d$ arbitrary $d$-blocks in between. When $M\ge0$, absorb the fixed prefix into the initialization and the fixed suffix into the terminal decoder; stochastic averaging keeps that decoder in $Y_j$. This produces an error-$\varepsilon$ clocked block machine for $G_\ell$, with registers at original cuts $p_\ell+dm$, $0\le m\le M$. Theorem~\ref{thm:radius54} bounds its narrow cuts $1\le m\le M$ by a constant depending only on $k$ and the fixed data. The omitted cuts of residue $\ell$ lie in a fixed prefix or a suffix of length less than $d$, and are uniformly bounded. If $M<0$, the whole horizon is bounded by $p_\ell+b_\ell$. Sum these bounds over the $d$ residue classes. This counts all cuts $1,\ldots,N$, independently of all other register widths.
\end{proof}
The constants in this theorem depend on the fixed alphabet, return period and certificate, interface, error, width $k$, and chosen finite prefixes. They do not depend on the horizon or on the machine. The original constants $B,d,\chi,L_k$ in Section~\ref{sec:local54} use $d$ for a contraction defect; the integer $d$ in this section is only a return period.

\begin{proposition}[Different phases can have different thresholds]\label{prop:parity55}
There is an alphabet with return period two for which $W_{N,0}=1$ on every even horizon, but $W_{N,\varepsilon}\to\infty$ along odd horizons for every $0\le\varepsilon<\rho/2$.
\end{proposition}
\begin{proof}
Take $H=(\R/2\pi\Z)\times\Z/2\Z$ and letters $(\theta,1),(-\theta,1)$, with $\theta/2\pi$ irrational. The generated group is dense in $H$. An exact return must have zero signed count and hence even length; the two-letter return gives every positive even length. Here $J=(\R/2\pi\Z)\times\{0\}$. For one binary-mean query put $F(\phi,0)=0$ and $F(\phi,1)=\rho\cos\phi$, where $0<\rho\le1$. The even-phase target is identically zero and needs one label. The odd-phase radius in the norm $|\cdot|/2$ is $\rho/2$. Apply Theorem~\ref{thm:phases55}.
\end{proof}

\section{Infinite component groups and boundary attainment}\label{sec:profinite55}
Let $H$ again be an arbitrary compact Hausdorff group, and set $K=H^\circ$. The quotient $H/K$ is profinite. We use two standard compact-group consequences of this fact and of Peter--Weyl theory \cite{Folland}: open normal subgroups form a neighborhood basis in a profinite group, and, for a continuous surjection $\pi:H\to G$ onto a compact Lie group, one has $\pi(K)=G^\circ$. To see the latter assertion, $\pi(K)$ is connected and normal, while $G/\pi(K)$ is a continuous group quotient of the profinite group $H/K$ and is therefore totally disconnected. This forces $G^\circ\subset\pi(K)$; the reverse inclusion is immediate. A continuous finite-dimensional representation of a profinite group has finite image, by the no-small-subgroups property of a Lie group and the open normal subgroup basis.

Define the component radius
\begin{equation}\label{eq:R055}
                 R_0=\max_{s,j,g\in H}\rad_{Y_j}(F_{s,j}(gK)).
\end{equation}
Unlike $r(U)$, this quantity need not itself be realized by any finite quotient.

\begin{lemma}[Approximation by finite quotient radii]\label{lem:infradius55}
With $r(U)$ as in \eqref{eq:quotientradius55},
\[
                 \inf_{U\text{ open normal in }H}r(U)=R_0.
\]
\end{lemma}
\begin{proof}
Every open subgroup contains $K$, so $r(U)\ge R_0$. Fix $\eta>0$. Uniform continuity, for the finite family of interfaces, gives an identity neighborhood $V$ such that
\[
 \norm{F_{s,j}(xv)-F_{s,j}(x)}_j<\eta
       \quad(x\in H,\ v\in V,\ s,j).
\]
There is an open normal subgroup $U$ containing $K$ and contained in $KV$, by the profinite quotient basis. Write each $u\in U$ as $kv$ with $k\in K$, $v\in V$. Then $F_{s,j}(gu)$ lies within $\eta$ of $F_{s,j}(gk)$. A center for the latter component image consequently covers the former coset image with an additional error at most $\eta$. Thus $r(U)\le R_0+\eta$.
\end{proof}

\begin{lemma}[A legal finite-dimensional interface approximation]\label{lem:Lieapprox55}
For every $\eta>0$ there is a continuous homomorphism $\pi:H\to G\subset O(D)$ onto a compact Lie group and continuous maps $\widetilde F_{s,j}:G\to Y_j$ such that
\[
           \max_{s,j}\norm{F_{s,j}-\widetilde F_{s,j}\circ\pi}_\infty<\eta.
\]
Moreover, for every $g\in H$,
\[
 \left|\rad_{Y_j}(F_{s,j}(gK))-
 \rad_{Y_j}(\widetilde F_{s,j}(\pi(g)G^\circ))\right|<\eta.
\]
\end{lemma}
\begin{proof}
Approximate all coordinates of all $F_{s,j}$ by finitely many real representative functions. The direct sum of their orthogonal representations is $\pi$, and the approximants are continuous functions on its compact image $G$. Choose Euclidean inner products on the finitely many output spaces. Project each approximant to the closed convex set $Y_j$ by nearest-point projection. This projection is continuous and nonexpansive in the chosen Euclidean norm. Norm equivalence permits the initial approximation to be chosen sufficiently accurate that the resulting $Y_j$-valued maps have error less than $\eta$ in the stated norms. The closed subgroup $G\subset O(D)$ is a compact Lie group. Since $\pi(K)=G^\circ$, the two indicated images have Hausdorff distance less than $\eta$. Their constrained radii therefore differ by less than $\eta$.
\end{proof}

\begin{theorem}[Arbitrary compact groups on both strict sides]\label{thm:profinite55}
For any compact-group interface, $\varepsilon>R_0$ implies a finite stationary permutation realization. At $\varepsilon=R_0$, such a realization exists exactly when $r(U)\le R_0$ for some open normal $U$.

Suppose in addition that executable exact returns have all sufficiently large lengths. If $\varepsilon<R_0$, then for each $k$ there is a horizon-independent upper bound on the number of cuts of width at most $k$, with all intervening widths unrestricted. In particular $W_{N,\varepsilon}\to\infty$. No finite-component assumption is made in either assertion.
\end{theorem}
\begin{proof}
The stationary assertions follow from Lemma~\ref{lem:infradius55} and Theorem~\ref{thm:finitequotient55}. For the lower bound choose a component image of radius greater than $\varepsilon$ and then $\eta>0$ with $2\eta$ smaller than that strict difference. Apply Lemma~\ref{lem:Lieapprox55}. A machine correct for the original interface at error $\varepsilon$ is correct for the approximating interface at error at most $\varepsilon+\eta$. The chosen component radius of that interface is greater than $\varepsilon+\eta$. The projected commands still have cofinite exact returns and generate $G$ densely. Theorem~\ref{thm:radius54} therefore gives the desired occupation bound for the very same registers. This argument never assumes that $K$ is open or that executable products lying exactly in $K$ are dense.
\end{proof}
At a positive boundary $R_0$, the stationary criterion just given is exact. We do not identify boundary attainment for all clocked families with that criterion. The zero-error boundary admits an additional finite-rank argument which does settle this identification.

\begin{theorem}[The exact clocked boundary]\label{thm:zero55}
Assume cofinite executable exact returns. The following are equivalent:
\begin{enumerate}
\item $\sup_NW_{N,0}<\infty$;
\item all $F_{s,j}$ factor through one finite continuous quotient $H/U$, with $U$ open normal;
\item there is an exact stationary finite permutation realization.
\end{enumerate}
This includes groups with infinitely many connected components.
\end{theorem}
\begin{proof}
The implications from the second assertion to the third and from the third to the first are immediate. Suppose the first holds with bound $k$. Theorem~\ref{thm:profinite55} implies $R_0=0$, so every interface is constant on each $K$-coset.

Fix a real coordinate $f$ of one $F_{s,j}$. For any finite families of executable prefix products $x_i$ and suffix products $y_l$, use cofinite returns to pad all selected prefix words to one common length $A$ and all suffix words to one common length $B$. Choose an exact realization at horizon $A+B$ of width at most $k$. At its cut $A$ the matrix
\[
                            (f(y_lx_i))_{i,l}
\]
factors through the $k$ hidden labels: the left factors are the actual state distributions after the padded prefixes and the right factors are the actual conditional terminal expectations on the padded suffixes. Hence this matrix has rank at most $k$. Padding does not need to preserve hidden states; it only preserves the physical products used in the displayed entries. By density and continuity the same rank bound holds for arbitrary $x_i,y_l\in H$.

It follows that the span of the left translates of $f$ in $C(H,\R)$ has dimension at most $k$. Otherwise $k+1$ linearly independent translates would have an invertible evaluation matrix at some $k+1$ points, contrary to the preceding rank bound. The elementary evaluation assertion follows inductively: a nonzero function is nonzero somewhere, and a function outside the span of the preceding ones cannot have all its evaluations forced by theirs.

Take the sum $V$ of these translate spaces over all seed--query coordinates. It is finite-dimensional and invariant under left translation. Translation is continuous in the uniform norm and preserves that norm, so it gives a continuous representation of $H$ on $V$ with compact image; Haar averaging makes it orthogonal. Because $K$ is normal and every original coordinate is constant on $K$-cosets, $K$ acts trivially on every translate and thus on $V$. The representation factors through $H/K$. Its image is finite by the profinite finite-dimensional representation fact above. Its kernel $U$ is open normal. Since every original coordinate belongs to $V$, all interfaces are $U$-invariant and descend to $H/U$.
\end{proof}

\begin{proposition}[A nonattained zero-radius boundary]\label{prop:adic55}
On a compact group with infinitely many components, every positive error can have a bounded stationary realization while the zero-error clocked width is unbounded.
\end{proposition}
\begin{proof}
Let $H=\Z_2$ be the additive group of $2$-adic integers, with commands $0$ and $1$. Their positive products are dense, and the zero command supplies exact returns of every length. If $x=\sum_{n\ge0}x_n2^n$, $x_n\in\{0,1\}$, prescribe the binary success probability
\[
                         f(x)=\sum_{n\ge0}\frac{2x_n}{3^{n+1}}.
\]
Uniform convergence makes $f$ continuous. It is injective: at the first differing digit the contribution $2/3^{n+1}$ exceeds the sum $1/3^{n+1}$ of all possible later differences. Since $H$ is totally disconnected, $K=\{0\}$ and $R_0=0$. Truncation after $m$ digits factors through $\Z/2^m\Z$ and has uniform error at most $3^{-m}$. Thus every positive tolerance admits a finite quotient machine. But the injective function $f$ cannot factor through any finite quotient. Theorem~\ref{thm:zero55} excludes bounded exact clocked width, and Theorem~\ref{thm:finitequotient55} excludes a finite exact stationary realization.
\end{proof}

\begin{theorem}[Linear exact width at a profinite boundary]\label{thm:adiclinear55}
For the interface of Proposition~\ref{prop:adic55},
\[
                         W_{N,0}=\Theta(N).
\]
More precisely, if $n$ is a power of two with $2n-1\le N$, then $W_{N,0}\ge n$, while $W_{N,0}\le N+1$. Every exact realization, with unrestricted intermediate widths, has at most $4k$ cuts in $\{1,\ldots,N\}$ of width at most $k$. Nevertheless, for each $\varepsilon>0$, $\sup_NW_{N,\varepsilon}<\infty$.
\end{theorem}
\begin{proof}
For nonnegative integers $z$ write $v_2(z+1)$ for the exponent of two in $z+1$. Incrementing the binary expansion, with $v_2(z+1)$ trailing ones, gives
\begin{equation}\label{eq:adicdifference55}
 b(z):=f(z+1)-f(z)=\frac53\,3^{-v_2(z+1)}-1.
\end{equation}
Let $n=2^m$, and form $B_n=(b(i+j))_{0\le i,j<n}$. Define a function on $\Z/n\Z$ by
\[
 \psi_m(x)=3^{-v_2(x)}\ (x\ne0),\qquad \psi_m(0)=3^{-m}.
\]
Here the valuation of a nonzero residue is independent of its representative. Since $1\le i+j+1\le2n-1$, its valuation equals that used in $\psi_m(i+j+1\bmod n)$, including the case $i+j+1=n$. Thus a permutation of the rows of $B_n$, namely $i\mapsto-i-1\pmod n$, makes it the circulant with kernel $(5/3)\psi_m-1$.

On $\Z/n\Z$ one has the elementary identity
\[
            \psi_m=1-2\sum_{\ell=1}^m3^{-\ell}
                               \mathbf1_{2^\ell\mid x}.
\]
The Fourier eigenvalue of this circulant at frequency zero is
\[
 \mu_0=\frac23 n6^{-m}>0.
\]
At frequency $q\ne0$ it is
\[
 \mu_q=-\frac{10n}{3}
           \sum_{\substack{1\le\ell\le m\\ n/2^\ell\mid q}}6^{-\ell}<0.
\]
Indeed a subgroup indicator has Fourier sum $n/2^\ell$ exactly when $n/2^\ell$ divides $q$, and zero otherwise. For $q\ne0$ the summand $\ell=m$ is always present. At $m=0$ the sole eigenvalue is $2/3$. Consequently $B_n$ is invertible for every dyadic $n$.

At an arbitrary command cut $t<N$, take a dyadic $n\le\min(t+1,N-t)$. For each $i<n$ choose a prefix of length $t$ with $i$ increment commands. For each $j<n$ choose two suffixes of length $N-t$ with respectively $j+1$ and $j$ increment commands. Fill all remaining positions with zero commands. The difference of their exact outputs is $b(i+j)$. Hence $B_n$ factors through the actual register $S_t$: one factor consists of the prefix state rows and the other of the differences of the two conditional suffix decoder columns. Invertibility gives $|S_t|\ge n$. The difference column need not be a legal decoder; it is used only for this rank bound on two legal experiments.

Taking $t=n-1$ proves $W_{N,0}\ge n$ whenever $2n-1\le N$. The largest dyadic $n\le(N+1)/2$ is greater than $(N+1)/4$, proving the linear lower order. Storing the number of increment commands uses at most $N+1$ labels and gives the exact upper bound.

For occupation, if $t<N$ and $|S_t|\le k$, the largest dyadic integer not exceeding $\min(t+1,N-t)$ is at most $k$, so $\min(t+1,N-t)<2k$. Such cuts lie within fewer than $2k$ positions of one of the two ends. Including the terminal cut gives at most $4k$ positive cuts. No width at another cut was restricted. The positive-error assertion is the finite-quotient truncation in Proposition~\ref{prop:adic55}.
\end{proof}
This supplies a sharp quantitative law on a genuinely infinite profinite group. Its linear exact-width obstruction and bounded positive-error realizations concern one fixed continuous interface; they are not a uniform approximation statement over all continuous outputs on that group.

\section{Binary, categorical, and observable interfaces}\label{sec:observable53}
\begin{corollary}[Component oscillation for binary means]\label{thm:classification53}
For a finite continuous binary-mean interface $f_{s,j}:H\to[-1,1]$, with the return hypothesis of Theorem~\ref{thm:radius54}, the exact bounded-width threshold in total variation is
\begin{equation}\label{eq:critical53}
 \epsc=\frac14\max_{s,j,c}\left(\max_{h\in c}f_{s,j}(h)-\min_{h\in c}f_{s,j}(h)\right).
\end{equation}
It is attained by a stationary permutation realization. Below it every fixed width has the occupation bound \eqref{eq:occupation54}.
\end{corollary}
\begin{proof}
In the norm $|\cdot|/2$, the radius of an interval is one quarter of its length. Its midpoint lies in $[-1,1]$. Apply Theorem~\ref{thm:radius54}.
\end{proof}
This recovers the scalar component theorem, including all of its nonuniform quantifiers. The vector theorem does not follow by applying this corollary separately to coordinates, because the coordinate centers need not belong to the output simplex.

\begin{proposition}[A three-outcome threshold]\label{prop:categorical54}
On the circle let $\omega=e^{2\pi i/3}$ and, for $j=0,1,2$, set
\[
 p_j(\phi)=\frac19\left|1+e^{i\phi}\omega^{-j}+e^{2i\phi}\omega^{-2j}\right|^2.
\]
For any dense circle command alphabet satisfying the return condition, the categorical law $p=(p_0,p_1,p_2)$ has critical total-variation error $2/3$. Each separately requested binary event $\{j\}$ has threshold $1/2$.
\end{proposition}
\begin{proof}
Orthogonality of the three characters gives $\sum_jp_j=1$; each entry is nonnegative. At $\phi=2\pi j/3$ the vector is the simplex vertex $e_j$. Hence every center $q\in\Delta_3$ has maximum distance at least
\[
 \max_j\TV(e_j,q)=1-\min_jq_j\ge2/3.
\]
The uniform center has distance at most $2/3$ from every point of $\Delta_3$, by convexity and the values at its vertices. Thus the constrained radius is exactly $2/3$. On the other hand each $p_j$ has range $[0,1]$, so its binary mean $2p_j-1$ has oscillation two. Apply Theorem~\ref{thm:radius54} and Corollary~\ref{thm:classification53}. In particular the categorical radius exceeds half the diameter, which here equals $1/2$.
\end{proof}

\begin{corollary}[Joint measure-once unitary laws]\label{cor:quantum54}
Let unitary commands have compact closure $H$ and satisfy the return condition. For a fixed unit seed $\psi_s$ and a finite measurement $E_{j,l}\ge0$, $\sum_lE_{j,l}=I$, prescribe the law
\[
 F_{s,j}(h)_l=\ip{h\psi_s}{E_{j,l}h\psi_s}.
\]
The critical error is its componentwise simplex radius in total variation. For closure $U(d)$ and one orthonormal-basis measurement, this radius is $1-1/d$.
\end{corollary}
\begin{proof}
The displayed map is a continuous simplex-valued polynomial in realified matrix entries. A closed unitary matrix group has finitely many components, so Theorem~\ref{thm:radius54} applies. For $U(d)$, every unit vector is reachable from the seed, and its squared coordinate magnitudes range over the whole simplex. The simplex vertices force radius at least $1-1/d$, attained by the uniform distribution.
\end{proof}
A query selects one complete measurement, not simultaneous answers to incompatible measurements. Several compatible finite outputs, with specified correlations, instead form one joint categorical law. Language recognition at a cutpoint is a different requirement.

For linear binary interfaces on a compact matrix group, write
\[
 f_{s,j}(h)=q_j^Thx_s\in[-1,1],\quad
 V_{\rm r}=\spanop\{hx_s:h\in H,s\in\mathcal S\},
\]
\[
 V_{\rm n}=\{x\in V_{\rm r}:q_j^Thx=0\text{ for all }h\in H,j\in\mathcal J\}.
\]
Both spaces are invariant. The observable space $V_{\rm o}=V_{\rm r}\cap V_{\rm n}^\perp$ is the orthogonal complement of $V_{\rm n}$ \emph{inside} $V_{\rm r}$, canonically realizing the observable quotient.

\begin{theorem}[Intrinsic zero-error criterion]\label{thm:observable53}
For these linear means, bounded exact clocked width is equivalent to finiteness of the generated group's image in $O(V_{\rm o})$, and to trivial action of $K$ on $V_{\rm o}$. An exact stationary permutation realization then exists.
\end{theorem}
\begin{proof}
Corollary~\ref{thm:classification53} identifies bounded exact width with constancy of all interface means on every component. Under that condition, for $k\in K$ and $g,h\in H$, normality gives
\[
 q_j^Th(k-I)gx_s=q_j^T(hkh^{-1})hgx_s-q_j^Thgx_s=0.
\]
Thus $(k-I)V_{\rm r}\subset V_{\rm n}$. On the invariant orthogonal complement $V_{\rm o}$ the difference also lies in $V_{\rm o}$, hence vanishes. The observable action factors through the finite component group. Conversely, if $K$ acts trivially on $V_{\rm o}$, the same difference is invisible and all means are componentwise constant. A finite image of the dense generated group has the same finite image closure from $H$, whose connected subgroup $K$ acts trivially. The zero-dimensional quotient is included.
\end{proof}

\begin{corollary}[The planar threshold]\label{cor:circle53}
Let an orthogonal planar alphabet contain $I$ and an irrational rotation, with any additional orthogonal letters. For signed coordinate seeds $\pm\rho e_1,\pm\rho e_2$, coordinate queries, and $0<\rho\le1$, the threshold is $\epsc=\rho/2$. At and above it one command label with a fair binary decoder suffices. Below it every fixed width has a finite occupation budget.
\end{corollary}
\begin{proof}
The identity component is $SO(2)$, and every specified mean has range $[-\rho,\rho]$ on each component of $H=SO(2)$ or $O(2)$. All midpoint decoders equal zero, so the seed/component register collapses to one label.
\end{proof}
For the noncommuting rational alphabet
\[
 I,\quad R=\begin{pmatrix}3/5&-4/5\\4/5&3/5\end{pmatrix},\quad
 F=\begin{pmatrix}1&0\\0&-1\end{pmatrix},
\]
one has $FRF=R^{-1}\ne R$. The rotation has infinite order: otherwise $(3+4i)/5$ and its inverse would be algebraic integers, whereas their rational sum $6/5$ is not an integer. This example has the sharp threshold above. We do not infer a badly approximable angle from rational matrix entries.

\begin{proposition}[Independent toral coordinates]\label{prop:torus53}
Suppose a component is identified with the full product torus $(\R/2\pi\Z)^m$, surjectively in independent coordinates, and a mean on it equals
\[
 a_0+\sum_{l=1}^m(a_l\cos\theta_l+b_l\sin\theta_l).
\]
Its contribution to \eqref{eq:critical53} is $\tfrac12\sum_l\sqrt{a_l^2+b_l^2}$.
\end{proposition}
\begin{proof}
Each summand has the symmetric range with amplitude $\sqrt{a_l^2+b_l^2}$, and all extrema are simultaneously attainable in the full product. Divide the resulting oscillation by four.
\end{proof}
On a proper subtorus, extrema must instead be taken subject to its relations. Ambient coordinates do not certify independence. Similarly, invisible directions do not create an obstruction: for $H=\{\diag(R_\theta,1)\}$, seed and query $e_3$ give constant mean one, whereas seed $\rho e_1$ and query $e_1$ have threshold $\rho/2$.

\begin{proposition}[Compact similarity]\label{prop:similarity53}
The classification applies to a finite alphabet in $GL(D,\R)$ whose generated group, including inverses, is uniformly bounded, with the executable return hypothesis retained. Conjugating to an invariant Euclidean metric changes neither output probabilities nor stochastic widths.
\end{proposition}
\begin{proof}
Uniform bounds on matrices and inverses make the closure a compact group. The matrix $Q=\int_Hh^Th\,d\mu(h)$ is positive definite and satisfies $A^TQA=Q$. Hence $Q^{1/2}AQ^{-1/2}$ is orthogonal. Transport the continuous output maps through this conjugacy; for a linear mean transport $x_s$ to $Q^{1/2}x_s$ and $q_j$ to $Q^{-1/2}q_j$. Exact identity products and their word lengths are preserved.
\end{proof}

\subsection{A curved categorical image with distinct component centers}
\begin{proposition}\label{prop:curved55}
Let $H=(\R/2\pi\Z)\times\Z/2\Z$, with an identity letter, an irrational rotation in the first factor, and a toggle of the second factor. Put
\[
 q_0=(\tfrac12,\tfrac13,\tfrac16),\qquad
 q_1=(\tfrac16,\tfrac13,\tfrac12),\qquad t=\tfrac1{12},
\]
and prescribe the three-outcome law
\[
 F(\theta,c)_l=q_{c,l}+t\cos(\theta-2\pi l/3),\qquad l=0,1,2.
\]
Each component image is a proper closed curve in the interior of the simplex, with unique constrained total-variation center $q_c$ and radius $t$. The union of both component images has radius $1/4$. Thus the sharp clocked threshold is $1/12$, and the stationary minimum at that threshold is exactly two labels.
\end{proposition}
\begin{proof}
The three cosines sum to zero and every coordinate is at least $1/12$, so $F$ is a legal law in the simplex interior. The first two independent trigonometric coordinates identify a nondegenerate ellipse in its affine plane. For a zero-sum vector $z\in\R^3$, one has $\frac12\norm{z}_1=\max_l|z_l|$. Each coordinate of the curve ranges over $[q_{c,l}-t,q_{c,l}+t]$. Hence for every center $y\in\Delta_3$,
\[
 \sup_\theta\TV(F(\theta,c),y)
       =t+\max_l|q_{c,l}-y_l|.
\]
This proves uniqueness of $q_c$ and radius $t$. For the union, the least possible value of $\max_{c,l}|q_{c,l}-y_l|$ is $1/6$, forced by the two endpoints in the first coordinate and attained by $y=(1/3,1/3,1/3)$. Its radius is therefore $t+1/6=1/4$. A one-label stationary machine has a constant decoder and cannot reach error $t$; two component labels attain it. The clocked threshold follows from Theorem~\ref{thm:radius54}.
\end{proof}

\begin{remark}
The total-variation diameter of a probability simplex is one, so half its diameter is $1/2$; the three-vertex example above has radius $2/3$. In the unitary corollary, $U(d)$ is connected. These facts explain respectively why coordinate midpoints do not solve the categorical problem and why the full-unitary component is the whole group.
\end{remark}

\section{Matching width laws throughout the subcritical interval}\label{sec:rates54}
We now distinguish the dimension of the command arithmetic from the dimension of the physical orbit. Let $s\ge1$ and let the alphabet contain the rotations $R_{2\pi\alpha_1},\ldots,R_{2\pi\alpha_s}$ and $I$. For the matching upper bound the alphabet is exactly these letters, optionally with the reflection $F=\diag(1,-1)$. Use signed coordinate seeds of radius $\rho$ and coordinate queries. Extra letters do not affect the lower bounds, but are not silently included in the upper construction.

For a real number $x$, write $\|x\|_{\R/\Z}$ for distance to the nearest integer. The dual Diophantine condition is
\begin{equation}\label{eq:DA54}
 \|n\cdot\alpha\|_{\R/\Z}\ge b\|n\|_1^{-\tau}
 \qquad(0\ne n\in\Z^s),\qquad b>0,\quad\tau\ge s.
\end{equation}
In particular $1,\alpha_1,\ldots,\alpha_s$ are rationally independent. When $\tau=s$ this is the dual badly approximable condition. All implicit constants below may depend on $\alpha,b,s,\rho,\varepsilon$, but not on $N$ or a stochastic realization.

\begin{theorem}[Full-subcritical matching law]\label{thm:matched54}
Suppose \eqref{eq:DA54} holds with $\tau=s$, and let $0<\rho<1$. For every fixed $0\le\varepsilon<\rho/2$,
\begin{equation}\label{eq:matched54}
                 W_{N,\varepsilon}=\Theta\bigl(N^{s/(2s+1)}\bigr).
\end{equation}
The upper bound is achieved by an exact horizon-dependent realization. At $\varepsilon\ge\rho/2$, one command label suffices. If $\tau\ge s$ is arbitrary, the lower bound $W_{N,\varepsilon}=\Omega(N^{s/(2\tau+1)})$ still holds throughout the same subcritical interval.
\end{theorem}
The exponent for one badly approximable rotation is $1/3$, not logarithmic. The statement includes zero error when the signal has strict amplitude slack $\rho<1$. The threshold theorem continues to include $\rho=1$, but the exact polynomial upper bound here is not asserted at that endpoint. The earlier arithmetic argument obtained this scale in a smaller noise range. The proof below supplies the missing full-noise lower estimate and states the exact upper realization in full.

\begin{proposition}[An explicit family in every command dimension]\label{prop:algebraicangles54}
Let $\xi$ be a real algebraic integer of degree $s+1$. Then
$\alpha=(\xi,\xi^2,\ldots,\xi^s)$ satisfies \eqref{eq:DA54} with $\tau=s$ and some effectively computable $b>0$.
\end{proposition}
\begin{proof}
For $0\ne n\in\Z^s$ put $H=\|n\|_1$, and choose $n_0\in\Z$ so that
$z=n_0+\sum_{l=1}^s n_l\xi^l$ has modulus $\|n\cdot\alpha\|_{\R/\Z}$.
The degree assumption gives $z\ne0$. Write $\xi_1=\xi,\ldots,\xi_{s+1}$ for the conjugates and put $M_j=\max_{1\le l\le s}|\xi_j|^l$. Then
$|n_0|\le H M_1+1/2$ and
\[
 \left|n_0+\sum_{l=1}^s n_l\xi_j^l\right|
 \le H(M_1+M_j+1/2)\qquad(j\ge2).
\]
The norm of the nonzero algebraic integer $z$ is a nonzero integer. Thus
\[
 |z|\ge H^{-s}\prod_{j=2}^{s+1}(M_1+M_j+1/2)^{-1}.
\]
Choose a positive rational $b$ below this computable positive constant and $1/2$. This proves the claim. It concerns algebraic \emph{angles}, not algebraic entries of their rotation matrices.
\end{proof}

\subsection{Harmonic localization}
For the harmonic calculation assume that the alphabet contains $I,R_\theta$ and that one seed-query pair has mean $\rho\cos\phi$ after accumulated angle $\phi$, where $0<\rho\le1$. Other seeds, queries, or commands impose additional constraints and do not invalidate the lower bound. Let $\lambda=e^{i\theta}$ have infinite order. Put $\delta=2\varepsilon<\rho$, and choose an integer
\begin{equation}\label{eq:r53}
 r\ge1,\qquad r>\frac{\delta}{\rho-\delta},\qquad m=r+1.
\end{equation}
For example $r=\lfloor\delta/(\rho-\delta)\rfloor+1$ works. Define
\begin{equation}\label{eq:circleconstants53}
 \Delta_r=2\left(\frac{\rho r}{r+1}-\delta\right)>0,\qquad
 A_r=\frac{2(1+\rho+\delta)(2r+1)m}{\Delta_r}>1.
\end{equation}

\begin{lemma}[Explicit harmonic localization]\label{lem:harmonic53}
The nonnegative trigonometric polynomials
\[
 p_\pm(\phi)=\frac{(1\pm\cos\phi)^r}{2^{-r}\binom{2r}{r}}
\]
have Haar mean one and weighted means
\[
 \int p_\pm(\phi)\rho\cos\phi\,\frac{d\phi}{2\pi}
                         =\pm\frac{\rho r}{r+1}.
\]
For the moving feature vector
\[
 Z_m(\phi)=(\cos\phi,\sin\phi,\ldots,\cos m\phi,\sin m\phi),
\]
any terminal realization with wordwise mean error at most $\delta$ obeys
\begin{equation}\label{eq:harmonicresidual53}
 \E\norm{\E[Z_m(\phi)\mid S]}\ge
 \frac{\Delta_r}{2(1+\rho+\delta)(2r+1)}.
\end{equation}
\end{lemma}
\begin{proof}
The identity $1+\cos\phi=2\cos^2(\phi/2)$, or the constant Laurent coefficient, gives
\[
 C_r:=\int(1+\cos\phi)^r\frac{d\phi}{2\pi}
            =2^{-r}\binom{2r}{r},\qquad
 C_{r+1}/C_r=\frac{2r+1}{r+1}.
\]
Hence the weighted cosine mean is $C_{r+1}/C_r-1=r/(r+1)$. Translation by $\pi$ gives the negative case.

We use the unweighted real coefficient vector in the basis $1,\cos\phi,\sin\phi,\ldots$, without orthonormal Fourier rescaling. The absolute sum of the real cosine coefficients of either $p_\pm$, including the constant coefficient, is
\[
 B_r=\frac{4^r}{\binom{2r}{r}}\le2r+1.
\]
For $p_+$ its coefficients are nonnegative and their sum is $p_+(0)$; for $p_-$ they differ only by alternating signs. The inequality follows because the central binomial coefficient is the largest of the $2r+1$ coefficients summing to $4^r$. Multiplication by $\rho\cos\phi$ has coefficient absolute sum at most $\rho B_r$, since $\cos u\cos v=(\cos(u+v)+\cos(u-v))/2$. Thus the Euclidean norms of the moving coefficient vectors of $p_\pm$ and $p_\pm\rho\cos\phi$ are at most $2r+1$ and $\rho(2r+1)$, respectively. For completeness, let $c=\E D$, $Q=\E\norm{\E[Z_m\mid S]}$. For either weight, the conditional coefficient identity gives $|c-\int p_\pm\rho\cos\phi|\le\delta+(1+\rho+\delta)(2r+1)Q$. Subtracting the two inequalities gives \eqref{eq:harmonicresidual53}. Correctness is averaged word by word before using this identity. The starting feature norm is $\sqrt m\le m$; $A_r$ deliberately uses the latter bound to obtain a rational closed form.
\end{proof}

The integer $r$ in \eqref{eq:r53} is the localization order; it is unrelated to the number $s$ of rotation letters. For fixed subcritical error it is finite and independent of $N$ and $k$.

\begin{theorem}[A full-noise arithmetic occupation bound]\label{thm:arithmetic54}
Assume \eqref{eq:DA54} and $0\le\varepsilon<\rho/2$. Fix the localization parameters $r,m,A_r$ in \eqref{eq:r53}--\eqref{eq:circleconstants53}. For an integer $k\ge1$, put
\[
 L=\left\lceil(2k)^{1/s}\right\rceil,\qquad B=s(L-1).
\]
Every error-$\varepsilon$ realization of peak at most $k$ satisfies
\begin{equation}\label{eq:finiteDA54}
 N<B\left(1+b^{-2}(mB)^{2\tau}\log A_r\right).
\end{equation}
With no restriction on other widths, the number of cuts of width at most $k$ is at most
\begin{equation}\label{eq:DAoccupation54}
 \left\lceil B-1+B b^{-2}(mB)^{2\tau}\log A_r\right\rceil.
\end{equation}
\end{theorem}
\begin{proof}
Use $s$ consecutive groups of $L-1$ command positions. In group $i$, choose $J_i$ uniformly from $\{0,\ldots,L-1\}$, apply $J_i$ copies of the $i$th rotation, and fill the other positions with identities. This defines an external law on length-$B$ words with $L^s\ge2k$ equally likely products. On the harmonic feature use
\[
 R_m(t)=\diag(R_t,R_{2t},\ldots,R_{mt}).
\]
For two distinct count vectors $J,J'$ and any frequency $1\le l\le m$, condition~\eqref{eq:DA54} and $\|l(J-J')\|_1\le mB$ give
\[
 |e^{2\pi i l(J-J')\cdot\alpha}-1|
 =2\bigl|\sin\bigl(\pi l(J-J')\cdot\alpha\bigr)\bigr|
 \ge4b(mB)^{-\tau}=:\sigma.
\]
Here $2\sin(\pi x)\ge4x$ for $0\le x\le1/2$. Also $b\le1/2$, $m\ge2$, and $B\ge1$, so $\sigma\le1$.

For a unit vector $u=(u_1,\ldots,u_m)$ in the frequency blocks, the squared distance between the two represented orbit points is
\[
 \sum_{l=1}^m\norm{u_l}^2
       |e^{2\pi i l(J-J')\cdot\alpha}-1|^2\ge\sigma^2.
\]
An open ball of radius $\sigma/2$ about any unit center contains at most one point. At least half the $L^s$ points therefore have distance at least $\sigma/2$ from all of any $k$ centers. Since the inner-product defect between unit vectors equals half their squared distance, the finite word law has defect at least
\[
                  d_k=\sigma^2/16=b^2(mB)^{-2\tau}.
\]
Apply conditional-centroid contraction to successive blocks and Lemma~\ref{lem:harmonic53} to the terminal query. The starting norm is $\sqrt m\le m$. Thus
\[
 \left\lfloor\frac NB\right\rfloor
 \le\frac{\log A_r}{-\log(1-d_k)}
 \le b^{-2}(mB)^{2\tau}\log A_r,
\]
which proves \eqref{eq:finiteDA54}. The greedy endpoint count, with identity gaps and no prefix, gives \eqref{eq:DAoccupation54}. Internal registers remain unrestricted. Finally $B<s(2k)^{1/s}$, so \eqref{eq:finiteDA54} implies the asserted $\Omega(N^{s/(2\tau+1)})$ lower bound.
\end{proof}

\begin{remark}
Testing a coordinate unit vector in \eqref{eq:DA54} gives $b\le1/2$. In the block construction, $L=\lceil(2k)^{1/s}\rceil\ge2$, hence $B=s(L-1)\ge1$. The integer $r$ is the localization order and $s$ is the number of rotations. For the algebraic-angle example, $1,\xi,\ldots,\xi^s$ are linearly independent over $\Q$ by the degree $s+1$ hypothesis; this ensures that the algebraic integer used in its norm bound is nonzero.
\end{remark}

\subsection{An exact polygon realization}
Write $\eta_\alpha(q)=\max_i\|q\alpha_i\|_{\R/\Z}$. The following elementary construction explains why a common denominator, rather than the number of numerical orbit points, controls the upper width.

\begin{proposition}[Exact enclosure upper bound]\label{prop:polygon54}
Let $0<\rho<1$, $a=(1+\rho)/2$, and $\gamma=\log(a/\rho)>0$. Choose an integer $q_0\ge4$ with $\cos(\pi/q_0)\ge a$. If $q\ge q_0$ and
\begin{equation}\label{eq:enclosure54}
                4\pi^2N\,\eta_\alpha(q)/q^2\le\gamma,
\end{equation}
then the indicated rotation alphabet, with or without $F$, has an exact horizon-$N$ realization with $q$ command labels.
\end{proposition}
\begin{proof}
Let $P_q$ be the regular polygon with vertices $v_j=(\cos(2\pi j/q),\sin(2\pi j/q))$, $0\le j<q$. For each $i$ choose a nearest integer $p_i$ to $q\alpha_i$ and put
\[
 t=\pi/q,\qquad h_i=2\pi|\alpha_i-p_i/q|\le t,\qquad
 \Lambda_i=\frac{\cos(t-h_i)}{\cos t}.
\]
Writing $n_j=(\cos((2j+1)t),\sin((2j+1)t))$, its facet description is
\[
               P_q=\{x:\ip{n_j}{x}\le\cos t\ (0\le j<q)\}.
\]
The facet inequalities of the regular polygon give $R_{2\pi\alpha_i}P_q\subset\Lambda_iP_q$: the nearest facet normal to a rotated vertex has angular difference $t-h_i$. Further,
\[
 \log\Lambda_i
 \le\log(1+\tan(t)h_i)\le\tan(t)h_i
 \le4\pi^2\eta_\alpha(q)/q^2,
\]
since $\cos h+\tan(t)\sin h\le1+\tan(t)h$ and $\tan t\le2t$ for $0\le t\le\pi/4$. Let $\Lambda=\max(1,\Lambda_1,\ldots,\Lambda_s)$. The standard reflection $F=\diag(1,-1)$ sends $v_j$ to $v_{-j\ ({\rm mod}\ q)}$ and preserves the chosen vertex set. The identity and reflection therefore preserve $P_q$, so every command $U$ satisfies $\Lambda^{-1}UP_q\subset P_q$.

For each command and each vertex choose a probability row with barycenter $\Lambda^{-1}Uv_j$ in $P_q$. These are actual common command rows, not independent prefix/suffix factorizations. At most three successors suffice by planar convexity. Initialize the seed $\rho u_s$, $u_s\in\{\pm e_1,\pm e_2\}$, with a row whose barycenter is $au_s$. This is possible since $a\le\cos(\pi/q)$, the polygon's inradius. After a length-$N$ word the mean vertex is
\[
                         a\Lambda^{-N}U_wu_s.
\]
By \eqref{eq:enclosure54}, $\rho\Lambda^N\le a$. Hence the terminal columns
\[
                       d_j(v)=\frac{\rho\Lambda^N}{a}\,e_j^Tv
\]
belong to $[-1,1]$ and return exactly $\rho e_j^TU_wu_s$. The rows can be stationary within a horizon; the terminal scaling and chosen denominator may depend on $N$. No one finite all-horizon exact machine is asserted.
\end{proof}

\begin{lemma}[Comparable simultaneous denominators]\label{lem:denominators54}
Assume \eqref{eq:DA54}. For every integer $Q\ge2$ there is $q$ with $1\le q\le Q^s$, $\eta_\alpha(q)\le1/Q$, and
\[
 q\ge c Q^{s/(\tau+1-s)},\qquad
 c=\left(\frac{b}{s^{\tau+1}}\right)^{s/(\tau+1-s)}.
\]
In particular, for $\tau=s$ one has $cQ^s\le q\le Q^s$.
\end{lemma}
\begin{proof}
Partition the torus into $Q^s$ half-open cubes and apply the pigeonhole principle to $0,\alpha,\ldots,Q^s\alpha$. Their difference gives $q\le Q^s$ and integers $p_i$ with $|q\alpha_i-p_i|\le1/Q$.

Put $M=\lfloor q^{1/s}\rfloor$. The $(M+1)^s>q$ vectors in $\{0,\ldots,M\}^s$ cannot all have distinct residues $n\cdot p$ modulo $q$, even when some nearest integers $p_i$ are negative. Subtract two to obtain a nonzero vector $n$ with $q\mid n\cdot p$ and $H:=\|n\|_1\le s q^{1/s}$. Therefore
\[
 bH^{-\tau}\le\|n\cdot\alpha\|_{\R/\Z}
 \le H/(qQ),\qquad
 q^{(\tau+1-s)/s}\ge bQ/s^{\tau+1}.
\]
The displayed bound follows. This argument supplies the needed quantitative transference directly; no unspecified denominator subsequence is being used.
\end{proof}

\begin{proof}[Proof of Theorem~\ref{thm:matched54}]
The lower estimate follows from Theorem~\ref{thm:arithmetic54}. For $\tau=s$, take the constants $a,\gamma,q_0,c$ above and an integer
\[
 Q\ge\max\left(2,(q_0/c)^{1/s},
                \left(\frac{4\pi^2N}{c^2\gamma}\right)^{1/(2s+1)}\right).
\]
Choosing the least such integer, Lemma~\ref{lem:denominators54} gives $q\le Q^s=O(N^{s/(2s+1)})$, $q\ge q_0$, and
\[
 4\pi^2N\eta_\alpha(q)/q^2
 \le4\pi^2N/(c^2Q^{2s+1})\le\gamma.
\]
Proposition~\ref{prop:polygon54} supplies the exact upper bound. At and above $\rho/2$ use the fair decoder from Corollary~\ref{cor:circle53}. This proves all claims.
\end{proof}
The command count $s$ is an arithmetic parameter, not a replacement for the dimension of the physical circle. The constants deteriorate as $\varepsilon\uparrow\rho/2$ through the localization degree, and as $\rho\uparrow1$ through the enclosure slack. Both dependences are explicit in the finite inequalities. The theorem makes no uniform claim at these limiting parameters.

\begin{corollary}[Unit amplitude at positive error]\label{cor:unitnoise55}
For the exact rotation alphabet of Theorem~\ref{thm:matched54}, with or without its specified reflection, suppose the angle vector is dual badly approximable. At $\rho=1$, for every fixed $0<\varepsilon<1/2$,
\[
                     W_{N,\varepsilon}=\Theta(N^{s/(2s+1)}).
\]
\end{corollary}
\begin{proof}
The harmonic localization and orbit-packing lower bound in Theorem~\ref{thm:arithmetic54} permits $\rho=1$ and every fixed $\varepsilon<1/2$. For the upper bound choose $0<\eta<2\varepsilon$ and put $\rho'=1-\eta\in(0,1)$, decreasing $\eta$ if necessary. Apply the exact polygon construction and comparable-denominator lemma to $\rho'$. On every seed, word, and query its mean differs from the unit-amplitude target by at most $\eta$, so the binary total-variation error is at most $\eta/2<\varepsilon$. Its width has the required order. This is a positive-error upper bound, not an exact realization at unit amplitude.
\end{proof}
The case $\rho=1$, $\varepsilon=0$ is not decided by this corollary. In particular, taking $\eta=0$ would destroy the strictly positive enclosure slack required by Proposition~\ref{prop:polygon54}.

\section{Effective algebraic-circle bounds}\label{sec:arithmetic53}
For algebraic unit-circle eigenvalues of infinite order, an explicit separation estimate replaces the Diophantine hypothesis of the preceding section. The parameters $r,m,A_r$ below are those of \eqref{eq:r53}--\eqref{eq:circleconstants53}; in this section $r$ denotes the localization order, not the number of command rotations. The alphabet contains $I,R_\theta$, and one interface mean is $\rho\cos\phi$.

\begin{lemma}[Separated powers give a block defect]\label{lem:separated53}
Suppose $|\lambda^n-1|\ge\sigma$ for every integer $1\le n\le2mk$, where $0<\sigma\le1$. Write $R_m(t)=\diag(R_t,R_{2t},\ldots,R_{mt})$ for the representation acting on $Z_m$. On this representation, the uniform law on the $2k$ words with products $R_\theta^j$, $0\le j<2k$, each padded to length $2k$, has defect at least $\sigma^2/16$ in \eqref{eq:gap54}.
\end{lemma}
\begin{proof}
Write a unit vector as $u=(u_1,\ldots,u_m)$ in its two-dimensional frequency blocks. For distinct indices $i,j\in\{0,\ldots,2k-1\}$,
\[
 \norm{R_m(i\theta)u-R_m(j\theta)u}^2
 =\sum_{l=1}^m\norm{u_l}^2|\lambda^{l(i-j)}-1|^2\ge\sigma^2.
\]
A ball of radius strictly less than $\sigma/2$ about any of $k$ centers contains at most one of these $2k$ points. Thus at least $k$ points have distance at least $\sigma/2$ from all centers (the same conclusion follows by taking the open balls of radius $\sigma/2$). For unit centers the defect is one half of the squared distance to the closest center. Averaging gives at least $(1/2)(\sigma^2/8)=\sigma^2/16$.
\end{proof}

\begin{theorem}[Full-noise algebraic-circle bound]\label{thm:algebraic53}
Assume $\lambda$ is algebraic and not a root of unity. Let its primitive irreducible integer polynomial be
\[
 P(z)=a z^d+a_{d-1}z^{d-1}+\cdots+a_0,\qquad a>0.
\]
Set $H_P=\max(a,|a_{d-1}|,\ldots,|a_0|)$ and
\[
                     B_P=a(1+H_P/a)^{d-1}>1.
\]
With $r,m,A_r$ as in \eqref{eq:r53}--\eqref{eq:circleconstants53}, every error-$\varepsilon$ realization of peak at most $k$ satisfies
\begin{equation}\label{eq:algwidth53}
 N<2k\left(1+16\,2^{2(d-1)}B_P^{4mk}\log A_r\right).
\end{equation}
Its number of command cuts $t\in\{1,\ldots,N\}$ with width at most $k$, even if other widths are unrestricted, is at most
\begin{equation}\label{eq:algoccupation53}
 2k-1+32k\,2^{2(d-1)}B_P^{4mk}\log A_r.
\end{equation}
Consequently, for fixed algebraic data, $\rho$, and $\varepsilon<\rho/2$,
\begin{equation}\label{eq:algrate53}
 W_{N,\varepsilon}\ge
 \frac{\log N-O(\log\log N)}{4m\log B_P}.
\end{equation}
All constants in the finite inequalities are explicit. If $\lambda=(a_1+ib_1)/c$ with integers $a_1,b_1$, $c>0$, $a_1^2+b_1^2=c^2$, and $\lambda$ not a root of unity, then \eqref{eq:algwidth53}--\eqref{eq:algrate53} hold with the factor $2^{2(d-1)}$ deleted and $B_P$ replaced by $c$.
\end{theorem}
\begin{proof}
List the conjugates with $\lambda_1=\lambda$. They are distinct by irreducibility in characteristic zero. None is a root of unity: otherwise irreducibility would make $P$ a cyclotomic polynomial up to its primitive sign. The resultant
\[
 \operatorname{Res}(P,z^n-1)=a^n\prod_{j=1}^d(\lambda_j^n-1)
\]
is therefore a nonzero integer. Each conjugate has modulus at most $R_P=1+H_P/a$. For completeness, if $|z|>R_P$, then
\[
 \sum_{j=0}^{d-1}|a_j||z|^j
 <\frac{H_P|z|^d}{|z|-1}<a|z|^d,
\]
so $P(z)\ne0$. Using $|z^n-1|\le2\max(1,|z|)^n$ for each other conjugate, the resultant bound gives
\begin{equation}\label{eq:resultant53}
 |\lambda^n-1|\ge
 2^{-(d-1)}\left(a\prod_{j=2}^d\max(1,|\lambda_j|)\right)^{-n}
 \ge 2^{-(d-1)}B_P^{-n}.
\end{equation}
For $1\le n\le2mk$, take $\sigma=2^{-(d-1)}B_P^{-2mk}$ in Lemma~\ref{lem:separated53}. The block defect is at least
\[
                  d_k=\frac{2^{-2(d-1)}B_P^{-4mk}}{16}.
\]
Repeated conditional contraction and Lemma~\ref{lem:harmonic53} imply
\[
 \left\lfloor\frac{N}{2k}\right\rfloor
 \le \frac{\log A_r}{-\log(1-d_k)}
 \le 16\,2^{2(d-1)}B_P^{4mk}\log A_r.
\]
This proves \eqref{eq:algwidth53}. The greedy occupation argument in the proof of Theorem~\ref{thm:radius54}, with $b=0$ and block length $2k$, proves \eqref{eq:algoccupation53}. Taking logarithms gives \eqref{eq:algrate53}: fix $C>1/(4m\log B_P)$. If $k\ge C\log N$ the claim is immediate; otherwise $\log k\le\log C+\log\log N$, which can be absorbed in the logarithm of \eqref{eq:algwidth53}. The implicit constant may depend on $P,\rho,\varepsilon$, not on $N$.

In the rational-circle case, $(a_1+ib_1)^n-c^n$ is a nonzero Gaussian integer, of modulus at least one. Thus $|\lambda^n-1|\ge c^{-n}$ directly. Repeat the proof with $\sigma=c^{-2mk}$. Here $c>1$ automatically. A nonminimal supplied denominator is still valid but weakens the bound.
\end{proof}
The threshold is sharp; these width rates are not asserted sharp. The factor $m$ records the increased localization degree as the error approaches the boundary. At smaller errors, the retained coordinate-calibration estimates can have better constants. One should use the larger of applicable lower bounds, rather than claiming numerical improvement at every parameter value.

\begin{proposition}[An explicit nonrational algebraic example]\label{prop:salem53}
Let $t=(1-\sqrt{13})/2$ and
\[
 \lambda=\frac{t+i\sqrt{4-t^2}}2.
\]
Then $|\lambda|=1$, $\lambda$ is not a root of unity, and one can take $d=4$ and $B_P=8$ in Theorem~\ref{thm:algebraic53}.
\end{proposition}
\begin{proof}
The number $t$ lies strictly between $-2$ and $2$ and solves $t^2-t-3=0$. Substitution of $t=z+z^{-1}$ gives
\[
                       P(z)=z^4-z^3-z^2-z+1.
\]
Its reduction modulo two is $z^4+z^3+z^2+z+1$. It has no linear factor over $\mathbb F_2$ and is not divisible by the only monic irreducible quadratic $z^2+z+1$; hence it is irreducible. Gauss's lemma proves irreducibility over $\mathbb Q$. The conjugate value $(1+\sqrt{13})/2>2$ gives a positive real conjugate of $\lambda$ greater than one. A root of unity cannot have such a conjugate. Finally $a=H_P=1$ and $d=4$, so $B_P=2^3=8$.
\end{proof}
For a reproducible rational calibration, take $\rho=1/10$ and $\varepsilon=1/25$. Then $\delta=2/25$, $r=5$, $m=6$, $\Delta_r=1/150$, and $A_r=23364$. Thus the rational-circle alphabet above obeys
\[
                 N<2k\bigl(1+16\cdot5^{24k}\log23364\bigr).
\]
This is a concrete bound inside the interval between the former small-error obstruction and the sharp threshold $1/20$. The exact coefficient and resultant checks accompanying the article verify the displayed arithmetic, not the universal theorem by finite experimentation.


\section{Finite encodings and boundary quotients}\label{sec:effective54}
The general continuous theorem is not an algorithm from arbitrary real data. We now give an explicit finite input class for which its conclusions are effective. A rational orthogonal input consists of a finite alphabet in $O(D,\Q)$ containing $I$, finite seed and query sets, and rational polynomials $F_{s,j,l}$ in the matrix entries specifying categorical probabilities. The promise $F_{s,j}(H)\subset\Delta_{M_j}$ can itself be checked once the compact closure is computed. Real algebraic output numbers are encoded by polynomial equations and isolating data.

We use two established algorithmic inputs: computation of a finite defining description of the Zariski closure of a finitely generated rational matrix group \cite{DJK,NPSW}, and real quantifier elimination with effective semialgebraic connected components \cite{BPR}. The former is a terminating closure algorithm, not merely an enumeration of valid polynomial identities with an unknown stopping point. Compact orthogonal groups are real algebraic, so the real points of the computed closure in $O(D)$ give the Euclidean closure $H$ \cite{BJKP}. The complex algebraic components must not be substituted for the real connected components used here.

\begin{theorem}[Effective rational classification]\label{thm:effective54}
For a rational orthogonal input, the following are computable:
\begin{enumerate}
\item a semialgebraic description of every connected component of $H$, the finite group $C$, and the component of each command;
\item the critical total-variation error $\epsc$ as an exact real algebraic number, and an attaining stationary component-permutation machine with algebraic decoders;
\item for each algebraic $0\le\varepsilon<\epsc$ and integer $k\ge1$, an integer occupation constant $L_k$ valid for every error-$\varepsilon$ clocked realization;
\item for every fixed $N$ and algebraic $\varepsilon\ge0$, the least width $W_{N,\varepsilon}$ and an algebraic realization attaining it.
\end{enumerate}
These are terminating computability statements. No polynomial running-time bound, practical implementation of general closure/quantifier elimination, or finite-random-bit cost bound is asserted.
\end{theorem}
\begin{proof}
Compute $H$ using the cited closure algorithm, then its real connected components and algebraic sample points. The component containing $I$ is $K$. Products of sample points and semialgebraic membership tests determine the component multiplication table and the command images. In particular $K$ is now an effectively given compact semialgebraic group.

For a component $c$ and a query, the condition that $q$ be a radius-$t$ center is the real first-order formula
\begin{equation}\label{eq:QEcenter54}
 q\in\Delta_{M_j},\qquad
 \forall h\in c:\quad
 \sum_{l=1}^{M_j}|F_{s,j,l}(h)-q_l|\le2t.
\end{equation}
Absolute values have a finite semialgebraic description. Quantifier elimination gives the set of admissible $(q,t)$. Its least $t$ exists by compactness and is real algebraic; algebraic sampling supplies an attaining $q$. Maximize over the finite interface and components. This proves (1)--(2), including equality at the critical error.

We detail effectivity of the occupation proof, since computing the threshold alone would not supply it. Choose a component with radius greater than the given $\varepsilon$, and enumerate positive words until one lies in that component. Its rational product gives the prefix used in the proof of Theorem~\ref{thm:radius54}. The resulting map $f(h)=F_{s,j}(hu_v)$ is polynomial, so it can be used directly in place of $G$, with no approximation error.

For any finite tensor degree, all required polynomial functions are coefficients of the explicitly given tensor representation. Its fixed subspace is
\[
 \{z:\ \forall h\in K,
                 (R(h)-I)z=0\}.
\]
This linear subspace is effectively semialgebraic. Real algebraic sampling, successively outside the span of the vectors already chosen, gives an algebraic basis in finitely many steps. Its Gram matrix gives the orthogonal projection $P$ exactly. Thus every polynomial Haar moment is computable algebraically from $P\Phi(e)$; numerical integration of Haar measure is unnecessary.

Enumerate finite lists of squares of rational polynomials on $K$, discarding zero Haar integrals and normalizing the others. For a list $p_i$, compute $b_i=\int p_if\,d\mu$ and its constrained simplex radius by quantifier elimination. Rational polynomials are uniformly dense among real polynomial restrictions on a compact matrix group; the localization proof therefore guarantees that some enumerated list has radius strictly greater than $\varepsilon$. Strict algebraic comparison detects that list. Fix a positive rational $\Delta$ below the radius difference and a positive rational upper bound $L$ for all coefficient constants in Proposition~\ref{prop:residual54}. These give the certified residual $Q\ge\Delta/L$.

Next enumerate finite rational probability laws on executable words of a common length whose products lie in $K$, together with rational $d\in(0,1)$. The assertion \eqref{eq:gap54}, quantified over the finitely many unit vectors in the algebraic moving space, is semialgebraic: each maximum of finitely many inner products can be expressed by a finite disjunction. It is therefore decidable. Lemma~\ref{lem:gap54} and sufficiently small rational perturbations of its atom weights ensure that this search terminates. Padding by the identity letter makes every finite choice executable at one common length.

The resulting $B,d,L,\Delta,Q_0$ determine \eqref{eq:Lk54}. To avoid exact comparisons involving logarithms, take a rational $A'\ge\max(1,LQ_0/\Delta)$ and use
\[
 \frac{\log A}{-\log(1-d)}\le\frac{A'}d.
\]
An integer above $b+B-1+BA'/d$ is a computable occupation bound. This proves (3).

For (4), fix a candidate width $k$ and use $k$ labels at every cut, padding smaller registers with unused labels. Stochastic rows, initialization, and categorical decoder probabilities are finitely many variables in compact simplexes. For each of the finitely many length-$N$ words, the output vector is polynomial in these variables, and the target vector is rational. The total-variation inequalities are semialgebraic. Quantifier elimination therefore decides feasibility and supplies an algebraic point when feasible. A deterministic register storing the seed and entire prefix is a finite exact realization, so increasing $k$ must terminate. This proves the least-width claim for each finite horizon, without claiming a finite-horizon cutoff for all stationary or nonuniform realization questions.
\end{proof}
The distinction between an effective theorem and an implemented solver matters. The accompanying finite checks verify displayed exact examples and proof constants; they do not execute the general group-closure and quantifier-elimination construction. The finite-bit compiler for rational finite groups remains in the complete predecessor supplement with its separate hypotheses and costs. It is not used to reprice arbitrary algebraic decoders as cost-free finite-bit procedures.

\subsection{The least component-quotient realization}
The upper bound $|\mathcal S||C|$ can often be reduced. Here is an exact finite characterization of that reduction, with its scope stated explicitly. Put $X=\mathcal S\times C$. A \emph{component-quotient realization} has a deterministic state of the form $\pi(s,c)$, where $\pi:X\to Q$ is a surjection; its command update is induced by $(s,c)\mapsto(s,[u_a]c)$. It is stationary and has no additional state recording a representative word within a component.

Call a partition $\mathcal B$ of $X$ equivariant if
\[
 (s,c)\sim(s',c')\quad\Longrightarrow\quad
 (s,dc)\sim(s',dc')\qquad(d\in C).
\]
For a block $B\in\mathcal B$ and query $j$, set
\[
 A_{B,j}=\bigcup_{(s,c)\in B}F_{s,j}(c).
\]
\begin{theorem}[Finite quotient minimum]\label{thm:quotient54}
The least number of command labels in an error-$\varepsilon$ component-quotient realization is
\begin{equation}\label{eq:quotient54}
 \min\left\{|\mathcal B|:\mathcal B\text{ equivariant and }
         \rad_{Y_j}(A_{B,j})\le\varepsilon\text{ for every }B,j\right\},
\end{equation}
with value $+\infty$ if there is no admissible partition. For rational orthogonal categorical inputs this minimum and an attaining quotient are computable. If $\epsc=0$ and $\varepsilon=0$, the coarsest admissible partition is the future-output equivalence
\[
 (s,c)\sim(s',c')\quad\Longleftrightarrow\quad
 F_{s,j}(dc)=F_{s',j}(dc')\quad\text{for every }d\in C,j,
\]
where componentwise constant maps are identified with their values.
\end{theorem}
\begin{proof}
An equivariant partition induces well-defined command permutations on its blocks. Choosing a minimizing center for each $A_{B,j}$ gives exactly the stated error. Conversely, the fibers of $\pi$ in any component-quotient realization form an equivariant partition. Its decoder must approximate every executable product in each corresponding component. Such products are dense in that component, so continuity extends the error bound to $A_{B,j}$. Its decoder is therefore a permitted center. This proves \eqref{eq:quotient54}. The command generators suffice to check equivariance, since their images generate the finite group $C$.

There are finitely many partitions of $X$. For polynomial categorical data their radius tests are of the form \eqref{eq:QEcenter54}, with a finite disjunction over the relevant seed/component pairs, so the minimum is computable. In the exact componentwise constant case, equality of all future outputs is an equivariant equivalence relation. Every admissible partition refines it, and its own blocks have singleton output images for each query. It is therefore the unique coarsest admissible partition.
\end{proof}
The minimization in \eqref{eq:quotient54} is a deterministic quotient minimum, distinct from the unrestricted stationary minimum in Theorem~\ref{thm:stationaryminimum55} and from the unrestricted nonuniform minimum at the boundary. Randomized encodings can identify numerical behaviors without identifying component pairs. What is classified here is exactly the finite deterministic quotient construction that attains the radius threshold. Theorem~\ref{thm:effective54}(4) separately computes the unrestricted minimum at each fixed horizon.

\subsection{Stationary minimization by permutation closures}
\begin{theorem}[Effective unrestricted stationary minimum]\label{thm:effectivemin55}
For rational orthogonal commands generating $H$, rational polynomial categorical outputs, and algebraic $\varepsilon\ge0$, one can compute $M_\varepsilon$ and an attaining algebraic machine when it is finite. An identity letter is not required. At each proposed width $k$, at most $(k!)^{|\mathcal A|}$ augmented orthogonal group closures and finite real quantifier-elimination problems suffice.
\end{theorem}
\begin{proof}
First compute $H$, its real connected components, and their constrained radii as in Theorem~\ref{thm:effective54}(1)--(2); those steps do not use an identity letter. The group $H\subset O(D)$ has finitely many components. Theorem~\ref{thm:finitequotient55} shows that $M_\varepsilon=\infty$ below their maximum radius. Otherwise $M_\varepsilon\le |\mathcal S||H/H^\circ|$.

For $k$ between one and this upper bound enumerate all command-permutation tuples $\sigma$. The matrices
\[
                         \diag(u_a^{-1},\Pi_a)\in O(D+k,\Q)
\]
generate a compact closure which is exactly the joint group $\mathcal L_\sigma$ of Theorem~\ref{thm:stationaryminimum55}. Compute it by the same algebraic closure procedure. In \eqref{eq:permutationfeasibility55}, the unknown seed rows and decoder rows are constrained to simplexes. The mean output is bilinear in these rows, the group variable is restricted by the computed polynomial equations, and $h^{-1}=h^T$ in the physical block. Consequently all target coordinates are polynomial in the universally quantified matrix entries. Total variation is semialgebraic. Real quantifier elimination decides the resulting existential--universal formula and supplies algebraic rows whenever it is feasible. The first feasible $k$ is the minimum by purification. The finite upper bound proves termination of the entire search.
\end{proof}
This is a structural reduction of stationary minimization, not merely an enumeration of arbitrary stochastic transition parameters. It does not equate stationary and horizon-specific minima. The two-cycle/three-cycle example in Proposition~\ref{prop:C655} already prevents replacing the permutation-feasibility test by deterministic component partitions.

\subsection{Algorithmic cost and arithmetic representation}
There are three distinct layers of effectivity. First, the group-closure procedure produces explicit defining equations; real-component decomposition and algebraic sampling then identify the component table. Second, radius computation and the stationary permutation tests are finite quantifier-elimination tasks on those equations. At fixed $k$, the stationary seed and output variables number $k(|\mathcal S|+\sum_jM_j)$ before normalization equalities are eliminated; the group variables can be taken among the $(D+k)^2$ augmented matrix entries. Thus the displayed finite search is independent of a horizon parameter. Its factorial number of discrete cases and the cost of the real-algebraic subproblems should not be mistaken for an efficient algorithm. No bound for the full closure-to-machine process in terms of only the original bit length is extracted here.

Third, occupation certificates in Theorem~\ref{thm:effective54}(3) add searches over localizer degrees and executable word laws. Their termination follows from strict analytic margins; the present proof supplies no a priori useful bound on the successful degree or word length. Fixed-horizon minimization, by contrast, has finitely many word constraints at each $N$ but their number is exponential in $N$ before repetitions are identified. These stages should be priced separately rather than hidden behind the word ``effective.''

All algebraic Haar moments in the localizer search are represented in a real algebraic number field with isolating data, using the fixed-subspace projector and its algebraic Gram inverse. Signs and strict comparisons are exact real-algebraic comparisons; a quadrature tolerance is not a certificate. Every rational command word has a rational matrix product and hence a rational target vector for rational polynomial readouts. The orthogonal closure fact used above is precisely that the Euclidean closure of a subgroup of $O(D)$ is the real zero set, inside $O(D)$, of its vanishing polynomial ideal, as used in \cite{BJKP}; replacing real connected components by complex algebraic ones would give the wrong quotient.

\section{Relation to numerical automata and positive realization}\label{sec:comparison54}
The input/output object and the order of quantifiers are decisive in comparing realization theorems. Classical probabilistic automata assign numerical acceptance probabilities to words, while stochastic-language results also study threshold recognition \cite{Rabin,Paz}. A numerical equivalence of two fixed machines is not the assertion that for every horizon there exists a new bounded-width approximate machine. Conversely, language equivalence need not preserve the numerical probabilities.

For weighted automata over a field, the finite-Hankel-rank characterization gives finite linear representations and identifies their minimum dimension \cite{BR}. It does not enforce nonnegative rows, stochastic normalization, or simplex-valued decoders. Our common-row contraction concerns these constraints simultaneously. The invariant observable quotient in Theorem~\ref{thm:observable53} uses the same reachable/invisible linear reduction; its additional conclusion is a finite-group criterion for exact positive realization with the full horizon-dependent quantifier.

The classical positive-realization problem asks for a positive system representing specified input/output data, and distinguishes positive dimension from unrestricted linear dimension \cite{BF04}. Invariant cones and polyhedral sections are part of that theory, not new notions introduced here. Proposition~\ref{prop:polygon54} is an explicit common invariant-enclosure construction. The new matched statement pairs it with a full-subcritical lower estimate against all clocked stochastic machines. Neither one cutwise nonnegative factorization nor a finite-rank Hankel matrix supplies that lower estimate.

Brodsky and Pippenger identify bounded-error measure-once quantum languages as group languages \cite[Theorem 3.3]{BP}. Their probabilistic simulation result preserves a cutpoint language with a potentially changed cutpoint \cite[Theorem 3.6]{BP}; it is not uniform additive preservation of every acceptance probability. Our categorical unitary corollary fixes an entire terminal measurement law in total variation and allows horizon-dependent stochastic approximants. The distinction is already visible in the $d$-outcome threshold $1-1/d$.

The closure theorem and strict-threshold decision method of Blondel, Jeandel, Koiran, and Portier \cite{BJKP}, and the effective Zariski-closure algorithms \cite{DJK,NPSW}, are algorithmic premises of Section~\ref{sec:effective54}. Computing the closure is not a new claim of this article. The additional outputs are a constrained minimax noise threshold, an attaining categorical decoder, and a terminating construction of an all-machine occupation constant. Recent Lie-algebra-based closure computations \cite{deGraaf} give an alternative representation of the algebraic closure, but their complex connected components still require real-component analysis for the present application.

Peter--Weyl approximation, Haar averaging, and the representative-function algebra are classical compact-group and almost-periodic-function machinery \cite{Folland}. Uniform approximation by a finite sum of matrix coefficients gives a finite \emph{linear} representation of an output function. It does not give positive stochastic rows or the repeated information loss quantified by \eqref{eq:occupation54}. Similarly, Steinberg's minimal compact topological automaton and enveloping-monoid results \cite[Theorems 2.2 and 2.6]{Steinberg} concern recognition of a language. His compact topological-monoid criterion \cite[Proposition 2.8]{Steinberg} uses clopen recognition, not a continuous simplex-valued response with a prescribed error norm. Our return condition does not classify arbitrary irreversible semigroups.

Distribution-based bisimulation metrics provide robust comparisons for specified probabilistic automata \cite{FSZ}. They do not by themselves impose a minimum over all horizon-specific stochastic realizations of an external group experiment. The finite future-output equivalence in Theorem~\ref{thm:quotient54} is deliberately a component-quotient statement, rather than a claim that approximate equivalence always has a unique minimal positive stochastic realization.

\begin{center}\small
\begin{tabular}{p{.27\textwidth}p{.31\textwidth}p{.32\textwidth}}
\toprule
Comparison object & Error or equivalence & Machine quantifier\\
\midrule
Weighted series \cite{BR} & Exact scalar coefficients; linear rank & One linear representation\\
Group-language simulation \cite{BP} & Cutpoint language; not additive probability error & One automaton\\
Positive realization \cite{BF04} & Exact specified numerical response & One positive system\\
Topological recognition \cite{Steinberg} & Clopen language acceptance & One recognizing action\\
Bisimulation metrics \cite{FSZ} & Distributional comparison of given systems & Two specified automata\\
Theorem~\ref{thm:radius54} & Worst-word norm error; TV for a whole categorical law & A new clocked stochastic machine at every horizon\\
Theorem~\ref{thm:matched54} & Binary TV, every fixed $\varepsilon<\rho/2$ & Same nonuniform lower model; exact upper construction\\
\bottomrule
\end{tabular}
\end{center}
The contribution combines stationary purification and its exact finite-quotient criterion with a sharp constrained clocked radius, phasewise arbitrary-width occupation, and, for the specified Diophantine family, a matching full-subcritical growth law. The elementary radius center, the arithmetic pigeonhole argument, the closure algorithm, and the harmonic-analysis ingredients are not asserted individually new. The accompanying theorem-level audit records the closest comparisons and their changed hypotheses; it is not a claim of exhaustive priority clearance.


\subsection{The classical matrix-semigroup boundary}
Flor's description of nonnegative matrix groups \cite[Theorems 2 and 3]{Flor} identifies their idempotent geometry and the finiteness of bounded groups; his account explicitly credits Brown and Schwarz for earlier compact and stochastic cases. We do not claim this matrix structure as new. Theorem~\ref{thm:purification55} uses a minimal-rank idempotent in a joint semigroup, not just in the matrix projection: its physical coordinate must be the group identity. Sandwiching each command then preserves all uniform error inequalities by closure, even when the original stochastic rows violate every nontrivial physical group relation. The resulting initialization may be random, a point essential to exact state minimization.

The stationary theorem has no return hypothesis. The clocked theorem has a different quantifier and needs synchronization; the empty-return example and the phase-dependent example show why the distinction cannot be removed by terminology. The exact-return-length reduction is an additive-semigroup argument, whereas the profinite zero-error converse uses finite Hankel rank and continuous finite-dimensional representations. None of these statements classifies an arbitrary irreversible physical semigroup.

The finite-quotient existence theorem does not say that $[H:U]$ is the optimal number of stochastic labels. The actual minimum is the permutation-feasibility invariant in Theorem~\ref{thm:stationaryminimum55}; the finite cyclic example separates them. For infinitely many components, the infimum of finite-quotient radii can fail to be attained, even at zero. Finally, the sharp quantitative law remains a theorem for a specified Diophantine alphabet, not a general growth theorem for all noncommuting compact-group representations or the exact unit-amplitude endpoint. These distinctions delimit the proved statements without weakening any of the retained conclusions.

\subsection{Companion results}
The companion supplement gives additional results on compatible reachable sections, one-surplus geometry, stable sections, finite-horizon certificates, word distortion, two-state duality, finite-bit compilation, and global residual hierarchies. The scalar two-extreme argument and earlier quantitative constants are also retained there. The present main proofs are self-contained apart from the classical results explicitly cited; none of these companion derivations is used to conceal a missing step in a principal theorem.

\begin{thebibliography}{99}
\bibitem{BPR} S. Basu, R. Pollack, and M.-F. Roy, \emph{Algorithms in Real Algebraic Geometry}, second ed., Algorithms and Computation in Mathematics 10, Springer, Berlin, 2006.
\bibitem{BF04} L. Benvenuti and L. Farina, A tutorial on the positive realization problem, \emph{IEEE Transactions on Automatic Control} \textbf{49} (2004), no. 5, 651--664.
\bibitem{BR} J. Berstel and C. Reutenauer, \emph{Noncommutative Rational Series with Applications}, Encyclopedia of Mathematics and its Applications 137, Cambridge University Press, Cambridge, 2011.
\bibitem{BJKP} V. D. Blondel, E. Jeandel, P. Koiran, and N. Portier, Decidable and undecidable problems about quantum automata, \emph{SIAM Journal on Computing} \textbf{34} (2005), no. 6, 1464--1473.
\bibitem{BP} A. Brodsky and N. Pippenger, Characterizations of 1-way quantum finite automata, \emph{SIAM Journal on Computing} \textbf{31} (2002), no. 5, 1456--1478.
\bibitem{DJK} H. Derksen, E. Jeandel, and P. Koiran, Quantum automata and algebraic groups, \emph{Journal of Symbolic Computation} \textbf{39} (2005), nos. 3--4, 357--371.
\bibitem{FSZ} Y. Feng, L. Song, and L. Zhang, Distribution-based bisimulation and bisimulation metric in probabilistic automata, arXiv:1512.05027 (2015).
\bibitem{Flor} P. Flor, On groups of non-negative matrices, \emph{Compositio Mathematica} \textbf{21} (1969), no. 4, 376--382.
\bibitem{Folland} G. B. Folland, \emph{A Course in Abstract Harmonic Analysis}, second ed., Chapman \& Hall/CRC, Boca Raton, 2016.
\bibitem{deGraaf} W. A. de Graaf, Computing the Zariski closure of a finitely generated matrix group, arXiv:2505.02577v2 (2026).
\bibitem{Hall} B. C. Hall, \emph{Lie Groups, Lie Algebras, and Representations: An Elementary Introduction}, second ed., Graduate Texts in Mathematics 222, Springer, Cham, 2015.
\bibitem{NPSW} K. Nosan, A. Pouly, S. Schmitz, M. Shirmohammadi, and J. Worrell, On the computation of the Zariski closure of finitely generated groups of matrices, arXiv:2106.01853v6 (2025).
\bibitem{Paz} A. Paz, \emph{Introduction to Probabilistic Automata}, Academic Press, New York, 1971.
\bibitem{Rabin} M. O. Rabin, Probabilistic automata, \emph{Information and Control} \textbf{6} (1963), no. 3, 230--245.
\bibitem{Steinberg} B. Steinberg, Topological dynamics and recognition of languages, arXiv:1306.1468 (2013).
\end{thebibliography}

\end{document}
